\documentclass[11pt]{article}
\usepackage{Physical_AI_21_CFR_Part_312}
\addbibresource{Physical_AI_21_CFR_Part_312.bib}
\title{End-to-End Physical AI Oncology Clinical Trial Unification}
\author{Kevin Kawchak, ChemicalQDevice}
\date{17 March 2026}
\begin{document}

\begin{titlepage}
\centering
{\Large END-TO-END PHYSICAL AI ONCOLOGY CLINICAL TRIAL UNIFICATION\par}
\vspace{1.5cm}
{\Huge\bfseries Adaption: 21 CFR Part 312\par}
\vspace{0.5cm}
{\LARGE\bfseries Investigational New Drug Application\par}
\vspace{0.3cm}
{\Large Modified ECFR\par}
\vspace{0.8cm}
{\large Draft release\par}
\vspace{0.2cm}
{\large Released on 17 March 2026\par}
\vspace{0.2cm}
{\large \href{https://doi.org/10.5281/zenodo.19057628}{10.5281/zenodo.19057628}\par}
\vspace{0.2cm}
{\large CEO Kevin Kawchak, ChemicalQDevice\par}
\vfill
{\small The original 21 CFR Part 312 document is a work in the public domain under 17 U.S.C. \S~105, and may be used, reproduced, incorporated into other works, adapted, modified, translated or distributed under a public license. 21 CFR Part 312 --- Investigational New Drug Application. \href{https://www.ecfr.gov/current/title-21/chapter-I/subchapter-D/part-312}{Source ECFR}. This current work is not endorsed or sponsored by CFR. The original ECFR formatting was reconstructed into Markdown, and adapted further into LaTeX using Claude Code Opus 4.6.\par}
\end{titlepage}

\pagenumbering{roman}
\tableofcontents
\clearpage
\section*{Prefatory Note}

This adaptation extends the prior 21 CFR Part 312 --- Investigational New Drug Application to address the end-to-end integration of Physical AI systems, advanced robotics, digital twins, and AI/ML agents into oncology clinical trials. The prior 21 CFR Part 312 regulation established foundational procedures and requirements governing the use of investigational new drugs, including submission, review, and approval of INDs by the Food and Drug Administration. This current document adapts those procedures to encompass the specialized requirements of Physical AI oncology clinical trial unification, including autonomous surgical robots, therapeutic positioning systems, diagnostic needle-placement platforms, rehabilitative exoskeletons, companion monitoring systems, AI-driven decision support, simulation-validated procedures, and federated multi-site trial coordination.

The Unification Standard Level (USL) scoring framework (v1.4.0 through v1.8.0) provides quantitative readiness assessments for robotic platforms across four dimensions: simulation switching, AI integration, cross-robot sharing, and clinical trial collaboration. USL scores range from 1.0 to 10.0 and are referenced throughout this adaptation to establish minimum readiness thresholds for clinical deployment and to inform IND submission requirements for Physical AI components.

The national MCP-PAI standard (v1.2.0) defines the five-server Model Context Protocol topology (authorization, FHIR, DICOM, ledger, provenance), 23 standardized tools, six actor roles, and five cumulative conformance levels that govern data access, audit, and robot agent operations in Physical AI oncology trials. This protocol infrastructure underpins the data management, safety reporting, and electronic records requirements throughout this adaptation.

This adaptation should be read in conjunction with the physical-ai-oncology-trials repository (v2.5.0, DOI: 10.5281/zenodo.19057628), the national MCP-PAI oncology trials standard (v1.2.0, DOI: 10.5281/zenodo.18869776), the ICH E6(R3) Physical AI adaptation (DOI: 10.5281/zenodo.18973368), the 21 CFR Part 50 Physical AI adaptation (DOI: 10.5281/zenodo.19040707), and the federated learning framework (DOI: 10.5281/zenodo.18840880).

\section*{Document History}

\begin{description}[leftmargin=3.8cm,style=sameline]
\item[Prior Regulation] 21 CFR Part 312 --- Investigational New Drug Application. Originally published at 52 FR 8831, March 19, 1987, with subsequent amendments through 2023. The prior regulation established the baseline framework for IND submissions, clinical investigation phases, sponsor and investigator responsibilities, safety reporting, and expanded access provisions.
\item[Current Adaptation] End-to-End Physical AI Oncology Clinical Trial Unification: Adaption of 21 CFR Part 312, released 17 March 2026. This adaptation incorporates Physical AI requirements throughout for oncology clinical investigations involving autonomous robotic systems, digital twins, simulation frameworks, and AI/ML agents.
\item[Repository Refs] physical-ai-oncology-trials v2.5.0 (March 2026), national-mcp-pai-oncology-trials v1.2.0 (March 2026). These repositories provide validated code, tools, specifications, and documentation supporting this adaptation.
\end{description}

\clearpage
\pagenumbering{arabic}

\section*{Prior 21 CFR Part 312: Public Domain Notice}

This section preserves the public domain notice from the prior 21 CFR Part 312 regulation. The original 21 CFR Part 312 document is a work of the United States Government and is in the public domain under 17 U.S.C. \S~105. It may be used, reproduced, incorporated into other works, adapted, modified, translated or distributed under a public license. This current adaptation is not endorsed or sponsored by CFR. The original ECFR formatting was reconstructed into Markdown and adapted further into LaTeX. All adaptations, modifications, and additions for Physical AI are clearly identified throughout.

\textbf{Authority:} 21 U.S.C. 321, 331, 351, 352, 353, 355, 360bbb, 371; 42 U.S.C. 262.

\textbf{Source:} 52 FR 8831, Mar. 19, 1987, unless otherwise noted.

\textbf{Editorial Note:} Nomenclature changes to part 312 appear at 69 FR 13717, Mar. 24, 2004.

\clearpage

\section*{Change Summary Table}

The following table identifies every section of the prior 21 CFR Part 312 regulation that has been modified in this adaptation and summarizes the nature of the Physical AI additions. Sections not listed below retain their original text without modification.

\begin{longtable}{>{\raggedright\arraybackslash}p{2.2cm} >{\raggedright\arraybackslash}p{11.5cm}}
\toprule
\textbf{Section} & \textbf{Nature of Physical AI Additions} \\
\midrule
\endfirsthead
\toprule
\textbf{Section} & \textbf{Nature of Physical AI Additions} \\
\midrule
\endhead
\midrule
\multicolumn{2}{r}{\textit{Continued on next page}} \\
\bottomrule
\endfoot
\bottomrule
\endlastfoot
\S~312.1 & Extended scope to include Physical AI systems performing, assisting, or monitoring clinical procedures; added applicability to robotic platforms, simulation frameworks, and digital twins \\
\S~312.2 & Added Physical AI applicability provisions for IND exemptions; clarified that Physical AI system components require IND coverage when used in clinical investigations \\
\S~312.3 & Added 22 Physical AI definitions covering robot agents, USL, digital twins, simulation frameworks, MCP, conformance levels, hash-chained audit trails, federated learning, PCCP, and safety constructs \\
\S~312.6 & Extended labeling requirements to Physical AI system components and robot-administered investigational drugs \\
\S~312.7 & Extended promotion restrictions to Physical AI system capabilities associated with investigational drugs \\
\S~312.8 & Added Physical AI cost recovery provisions for robotic system deployment, simulation validation, and digital twin infrastructure \\
\S~312.10 & Extended waiver provisions to accommodate Physical AI system requirements \\
\S~312.20 & Added IND requirement for Physical AI system components used in clinical investigations \\
\S~312.21 & Added Physical AI considerations for each phase of clinical investigation, including simulation-based Phase 0 validation \\
\S~312.22 & Extended IND submission principles to Physical AI system documentation requirements \\
\S~312.23 & Added comprehensive Physical AI documentation requirements for IND content, including robot specifications, simulation validation data, digital twin models, AI/ML algorithms, and cybersecurity assessments \\
\S~312.30 & Extended protocol amendment requirements to Physical AI system modifications, algorithm updates, and simulation framework changes \\
\S~312.31 & Extended information amendment requirements to Physical AI technical updates \\
\S~312.32 & Added Physical AI safety reporting requirements including robot malfunction events, AI prediction errors, digital twin divergence, and cybersecurity incidents \\
\S~312.33 & Extended annual reporting to Physical AI system performance metrics, USL score updates, and simulation validation results \\
\S~312.38 & Extended IND withdrawal procedures to Physical AI system decommissioning \\
\S~312.40 & Added Physical AI general requirements for IND effectiveness, including simulation validation and USL threshold verification \\
\S~312.42 & Added Physical AI grounds for clinical hold including robotic safety failures, AI model degradation, and simulation-reality divergence \\
\S~312.44 & Extended termination grounds to Physical AI system failures and cybersecurity compromises \\
\S~312.45 & Extended inactive status provisions to Physical AI system dormancy and reactivation \\
\S~312.47 & Added Physical AI system demonstration requirements for regulatory meetings \\
\S~312.48 & Extended dispute resolution to Physical AI technical disagreements \\
\S~312.50 & Extended sponsor responsibilities to Physical AI system selection, validation, monitoring, and lifecycle management \\
\S~312.52 & Extended CRO transfer provisions to Physical AI obligations \\
\S~312.53 & Extended investigator and monitor selection to Physical AI qualifications \\
\S~312.54 & Extended emergency research monitoring to Physical AI systems \\
\S~312.55 & Extended investigator information requirements to Physical AI system updates \\
\S~312.56 & Extended ongoing investigation review to Physical AI system compliance monitoring \\
\S~312.57 & Extended recordkeeping to Physical AI system logs, simulation data, and digital twin records \\
\S~312.58 & Extended inspection access to Physical AI system records and audit trails \\
\S~312.59 & Extended drug disposition to Physical AI system decommissioning procedures \\
\S~312.60 & Extended investigator responsibilities to Physical AI system oversight \\
\S~312.61 & Extended drug control to Physical AI system access control \\
\S~312.62 & Extended investigator recordkeeping to Physical AI case histories and system interaction logs \\
\S~312.64 & Extended investigator reporting to Physical AI system events and performance metrics \\
\S~312.66 & Extended IRB review assurance to Physical AI system components \\
\S~312.68 & Extended inspection access to Physical AI investigator records \\
\S~312.69 & Extended controlled substance handling to Physical AI-mediated administration \\
\S~312.70 & Extended disqualification to Physical AI system misuse and noncompliance \\
\S~312.80 & Extended life-threatening illness provisions to Physical AI accelerated development pathways \\
\S~312.82 & Extended early consultation to Physical AI system design review \\
\S~312.83 & Extended treatment protocols to Physical AI-assisted therapeutic delivery \\
\S~312.84 & Extended risk-benefit analysis to Physical AI system risk profiles \\
\S~312.85 & Extended Phase 4 studies to Physical AI post-market performance monitoring \\
\S~312.87 & Extended active monitoring to Physical AI system clinical performance \\
\S~312.88 & Extended patient safety safeguards to Physical AI-specific protections \\
\S~312.110 & Extended import/export requirements to Physical AI system components and software \\
\S~312.120 & Extended foreign clinical study acceptance to Physical AI system validation data \\
\S~312.130 & Extended public disclosure provisions to Physical AI system data \\
\S~312.160 & Extended laboratory use provisions to Physical AI preclinical testing \\
\S~312.300 & Extended expanded access scope to Physical AI-assisted treatment \\
\S~312.305 & Extended expanded access requirements to Physical AI system safeguards \\
\S~312.310 & Extended individual patient access to Physical AI-assisted emergency treatment \\
\S~312.315 & Extended intermediate-size access to Physical AI multi-site deployment \\
\S~312.320 & Extended treatment IND to Physical AI widespread treatment use \\
New J & New Subpart J: Additional Requirements for Physical AI Investigational Systems in Oncology Trials \\
\end{longtable}

\clearpage

\section{SUBPART A --- GENERAL PROVISIONS}

\subsection{\S~312.1 Scope}

(a) This part contains procedures and requirements governing the use of investigational new drugs, including procedures and requirements for the submission to, and review by, the Food and Drug Administration of investigational new drug applications (IND's). An investigational new drug for which an IND is in effect in accordance with this part is exempt from the premarketing approval requirements that are otherwise applicable and may be shipped lawfully for the purpose of conducting clinical investigations of that drug.

(b) References in this part to regulations in the Code of Federal Regulations are to chapter I of title 21, unless otherwise noted.

\subsubsection{Physical AI Adaptation of \S~312.1}

(c) This part additionally applies to all clinical investigations in which Physical AI systems perform, assist, or monitor procedures involving the administration or dispensing of investigational new drugs to human subjects. Physical AI systems encompass the seven robot categories recognized for clinical investigation purposes: autonomous surgical robots, cobots, humanoid robots, therapeutic robots, diagnostic robots, assistive robots, and rehabilitative robots. These categories are further organized into five functional types defined in the national MCP-PAI standard (v1.2.0): autonomous surgical robots, therapeutic positioning systems, diagnostic needle-placement platforms, rehabilitative exoskeletons, and companion monitoring systems.

(d) The scope of this part extends to IND applications involving investigational drugs that are administered, dispensed, prepared, compounded, or monitored by Physical AI systems. When a Physical AI system participates in any step of the drug administration pathway, the IND must document the system's role, safety parameters, and operational specifications as described in \S~312.23 of this part. The processing of investigational drug data through Physical AI systems, including through the five MCP servers (trialmcp-authz, trialmcp-fhir, trialmcp-dicom, trialmcp-ledger, trialmcp-provenance), constitutes an activity subject to the requirements of this part.

(e) The Unification Standard Level (USL) framework serves as the readiness assessment methodology for Physical AI platforms used in clinical investigations under this part. USL scores are computed across four equally weighted dimensions (simulation switching, AI integration, cross-robot sharing, clinical trial collaboration) and range from 1.0 (Minimal) to 10.0 (Reference). The national MCP-PAI standard specifies minimum USL thresholds by procedure type: surgical robots require USL 7.0 or above, therapeutic positioning and diagnostic needle placement require USL 6.0 or above, rehabilitative exoskeletons require USL 4.0 or above, and companion monitoring systems require USL 3.0 or above. These thresholds apply to Physical AI systems used in clinical investigations conducted under INDs submitted pursuant to this part.

(f) Four simulation frameworks are recognized for pre-clinical validation of Physical AI systems before their use in IND-covered clinical investigations: NVIDIA Isaac Lab (v2.3.1) for GPU-accelerated robot training, MuJoCo (v3.4.0) for precision physics simulation, Gazebo Sim (v10.0.0) for ROS 2 integrated simulation, and PyBullet (v3.2.5) for rapid prototyping. Cross-framework validation via the unification bridge must demonstrate behavioral consistency with trajectory deviations below 2mm and force discrepancies below 0.5N before a Physical AI system may be used in clinical investigations under an IND.

(g) Digital twin models used for treatment planning, procedure rehearsal, or real-time intraoperative guidance in connection with investigational drug administration are subject to the documentation and validation requirements of this part. Digital twins must be calibrated to patient-specific imaging data and continuously validated against clinical outcomes throughout the investigation.

[52 FR 8831, Mar. 19, 1987; 46 FR 8979, Jan. 27, 1981. Physical AI adaptation added 17 March 2026.]

\subsection{\S~312.2 Applicability}

(a) \textit{Applicability.} Except as provided in this section, this part applies to all clinical investigations of products that are subject to section 505 of the Federal Food, Drug, and Cosmetic Act or to the licensing provisions of the Public Health Service Act (58 Stat. 632, as amended (42 U.S.C. 201 \textit{et seq.})).

(b) \textit{Exemptions.}

\begin{description}[leftmargin=0.5cm,style=sameline]
\item[(1)] The clinical investigation of a drug product that is lawfully marketed in the United States is exempt from the requirements of this part if all the following apply:

\begin{description}[leftmargin=0.8cm,style=sameline]
\item[(i)] The investigation is not intended to be reported to FDA as a well-controlled study in support of a new indication for use nor intended to be used to support any other significant change in the labeling for the drug;
\item[(ii)] If the drug that is undergoing investigation is lawfully marketed as a prescription drug product, the investigation is not intended to support a significant change in the advertising for the product;
\item[(iii)] The investigation does not involve a route of administration or dosage level or use in a patient population or other factor that significantly increases the risks (or decreases the acceptability of the risks) associated with the use of the drug product;
\item[(iv)] The investigation is conducted in compliance with the requirements for institutional review set forth in part 56 and with the requirements for informed consent set forth in part 50; and
\item[(v)] The investigation is conducted in compliance with the requirements of \S~312.7.
\end{description}

\item[(2)] \hfill
\begin{description}[leftmargin=0.8cm,style=sameline]
\item[(i)] A clinical investigation involving an in vitro diagnostic biological product listed in paragraph (b)(2)(ii) of this section is exempt from the requirements of this part if (\textit{a}) it is intended to be used in a diagnostic procedure that confirms the diagnosis made by another, medically established, diagnostic product or procedure and (\textit{b}) it is shipped in compliance with \S~312.160.
\item[(ii)] In accordance with paragraph (b)(2)(i) of this section, the following products are exempt from the requirements of this part: (\textit{a}) blood grouping serum; (\textit{b}) reagent red blood cells; and (\textit{c}) anti-human globulin.
\end{description}

\item[(3)] A drug intended solely for tests in vitro or in laboratory research animals is exempt from the requirements of this part if shipped in accordance with \S~312.160.
\item[(4)] FDA will not accept an application for an investigation that is exempt under the provisions of paragraph (b)(1) of this section.
\item[(5)] A clinical investigation involving use of a placebo is exempt from the requirements of this part if the investigation does not otherwise require submission of an IND.
\item[(6)] A clinical investigation involving an exception from informed consent under \S~50.24 of this chapter is not exempt from the requirements of this part.
\end{description}

(c) \textit{Bioavailability studies.} The applicability of this part to in vivo bioavailability studies in humans is subject to the provisions of \S~320.31.

(d) \textit{Unlabeled indication.} This part does not apply to the use in the practice of medicine for an unlabeled indication of a new drug product approved under part 314 or of a licensed biological product.

(e) \textit{Guidance.} FDA may, on its own initiative, issue guidance on the applicability of this part to particular investigational uses of drugs. On request, FDA will advise on the applicability of this part to a planned clinical investigation.

\subsubsection{Physical AI Adaptation of \S~312.2}

(f) \textit{Physical AI system components.} When a clinical investigation involves both an investigational new drug and a Physical AI system that administers, monitors, or otherwise participates in the drug delivery pathway, the IND requirements of this part apply to the combined drug-device-robot system. The exemptions in paragraph (b)(1) of this section do not apply if the introduction of a Physical AI system into the drug administration pathway significantly increases the risks or introduces new categories of risk (including robotic malfunction, AI prediction error, cyber-physical failure, or human-robot interaction hazards) not present in conventional drug administration.

(g) \textit{Simulation-only investigations.} Investigations conducted entirely in simulation environments (NVIDIA Isaac Lab, MuJoCo, Gazebo Sim, or PyBullet) without any human subject involvement are not subject to the requirements of this part. However, simulation data generated for the purpose of supporting a subsequent IND submission must comply with the documentation standards specified in \S~312.23 and must be generated using validated simulation frameworks meeting the specifications in the physical-ai-oncology-trials repository (v2.5.0).

(h) \textit{Digital twin investigations.} Clinical investigations that use patient-specific digital twin models for treatment planning or procedure rehearsal, where the digital twin output directly influences the administration of an investigational drug to a human subject, are subject to the requirements of this part. Digital twin models used solely for post-hoc analysis or research purposes that do not influence active treatment decisions are not subject to IND requirements but must be documented in the annual report under \S~312.33.

[52 FR 8831, Mar. 19, 1987, as amended at 61 FR 51529, Oct. 2, 1996; 64 FR 401, Jan. 5, 1999. Physical AI adaptation added 17 March 2026.]

\subsection{\S~312.3 Definitions and Interpretations}

(a) The definitions and interpretations of terms contained in section 201 of the Act apply to those terms when used in this part:

(b) The following definitions of terms also apply to this part:

\textit{Act} means the Federal Food, Drug, and Cosmetic Act (secs. 201--902, 52 Stat. 1040 \textit{et seq.}, as amended (21 U.S.C. 301--392)).

\textit{Clinical investigation} means any experiment in which a drug is administered or dispensed to, or used involving, one or more human subjects. For the purposes of this part, an experiment is any use of a drug except for the use of a marketed drug in the course of medical practice.

\textit{Contract research organization} means a person that assumes, as an independent contractor with the sponsor, one or more of the obligations of a sponsor, e.g., design of a protocol, selection or monitoring of investigations, evaluation of reports, and preparation of materials to be submitted to the Food and Drug Administration.

\textit{FDA} means the Food and Drug Administration.

\textit{IND} means an investigational new drug application. For purposes of this part, ``IND'' is synonymous with ``Notice of Claimed Investigational Exemption for a New Drug.''

\textit{Independent ethics committee (IEC)} means a review panel that is responsible for ensuring the protection of the rights, safety, and well-being of human subjects involved in a clinical investigation and is adequately constituted to provide assurance of that protection. An institutional review board (IRB), as defined in \S~56.102(g) of this chapter and subject to the requirements of part 56 of this chapter, is one type of IEC.

\textit{Investigational new drug} means a new drug or biological drug that is used in a clinical investigation. The term also includes a biological product that is used in vitro for diagnostic purposes. The terms ``investigational drug'' and ``investigational new drug'' are deemed to be synonymous for purposes of this part.

\textit{Investigator} means an individual who actually conducts a clinical investigation (\textit{i.e.}, under whose immediate direction the drug is administered or dispensed to a subject). In the event an investigation is conducted by a team of individuals, the investigator is the responsible leader of the team. ``Subinvestigator'' includes any other individual member of that team.

\textit{Marketing application} means an application for a new drug submitted under section 505(b) of the act or a biologics license application for a biological product submitted under the Public Health Service Act.

\textit{Sponsor} means a person who takes responsibility for and initiates a clinical investigation. The sponsor may be an individual or pharmaceutical company, governmental agency, academic institution, private organization, or other organization. The sponsor does not actually conduct the investigation unless the sponsor is a sponsor-investigator. A person other than an individual that uses one or more of its own employees to conduct an investigation that it has initiated is a sponsor, not a sponsor-investigator, and the employees are investigators.

\textit{Sponsor-Investigator} means an individual who both initiates and conducts an investigation, and under whose immediate direction the investigational drug is administered or dispensed. The term does not include any person other than an individual. The requirements applicable to a sponsor-investigator under this part include both those applicable to an investigator and a sponsor.

\textit{Subject} means a human who participates in an investigation, either as a recipient of the investigational new drug or as a control. A subject may be a healthy human or a patient with a disease.

[52 FR 8831, Mar. 19, 1987, as amended at 64 FR 401, Jan. 5, 1999; 64 FR 56449, Oct. 20, 1999; 73 FR 22815, Apr. 28, 2008]

\subsubsection{Physical AI Definitions}

The following additional definitions apply to clinical investigations involving Physical AI systems. These definitions are derived from the physical-ai-oncology-trials repository (v2.5.0) and the national MCP-PAI oncology trials standard (v1.2.0).

(c) \textit{Physical AI system} means an AI system operating in the physical world through a robotic platform used in clinical investigations involving investigational new drugs. Physical AI systems include autonomous surgical robots, therapeutic positioning systems, diagnostic needle-placement platforms, rehabilitative exoskeletons, and companion monitoring systems. The term encompasses the robotic hardware, its control software, embedded AI/ML models, and associated sensor systems that collectively perform, assist, or monitor clinical procedures involving the administration or dispensing of investigational drugs to human subjects.

(d) \textit{Robot agent} means an autonomous Physical AI system executing clinical procedures within an oncology trial conducted under an IND, interacting through MCP-mediated data access. The robot agent operates under the six-step workflow defined in the national MCP-PAI standard: authenticate, retrieve clinical data, access imaging, execute procedure, record audit, and record provenance. Robot agents are one of six actor types (alongside Trial Coordinator, Data Monitor, Auditor, Sponsor, and CRO) in the MCP actor model.

(e) \textit{Unification Standard Level (USL)} means a scoring framework (v1.4.0 through v1.8.0) evaluating Physical AI platform readiness across four equally weighted dimensions: simulation switching, AI integration, cross-robot sharing, and clinical trial collaboration. USL scores range from 1.0 to 10.0, mapping to ten levels from Minimal (1.0) to Reference (10.0) and three readiness bands: Foundational (1.0 to 4.9), Intermediate (5.0 to 6.9), and Advanced (7.0 to 10.0). Nine robotic platforms have been evaluated: da Vinci dVRK (USL 7.1), Franka Emika Panda (USL 7.4), Boston Dynamics Atlas (USL 5.8), Kinova Gen3 (USL 5.7), Medtronic Hugo RAS (USL 4.5), Agility Digit (USL 4.2), Tesla Optimus (USL 3.6), CMR Versius (USL 3.4), and UFACTORY xArm 7 (USL 3.4).

(f) \textit{Digital twin} means a patient-specific computational model used for treatment simulation, procedure planning, and real-time intraoperative guidance in connection with the administration of an investigational drug. Digital twins are calibrated to longitudinal imaging data using reaction-diffusion equations, pharmacokinetic/pharmacodynamic models, or other mathematical frameworks and are continuously updated with real-world clinical data. Digital twin models used in IND-covered investigations must meet the validation requirements specified in this part.

(g) \textit{Simulation framework} means a physics-based computational environment used for validating Physical AI system behavior before clinical deployment in IND-covered investigations. Recognized simulation frameworks include NVIDIA Isaac Lab (v2.3.1), MuJoCo (v3.4.0), Gazebo Sim (v10.0.0), and PyBullet (v3.2.5). The validated integration path follows: training (Isaac Lab) to validation (MuJoCo) to hardware deployment to clinical use.

(h) \textit{Robot capability profile} means a machine-readable specification of a Physical AI platform's capabilities, safety prerequisites, USL score, and required MCP tools. The profile is defined by a JSON schema (robot-capability-profile.schema.json) in the national MCP-PAI standard and must be registered with the authorization server before any clinical procedure involving an investigational drug.

(i) \textit{Task order} means a scheduled clinical trial task assigned to a Physical AI system, including procedure type, robot assignment, safety check requirements, and lifecycle state. Task orders progress through defined states: scheduled, safety\_check, in\_progress, and completed, cancelled, or failed.

(j) \textit{Emergency stop} means an immediate halt capability that must be available for any Physical AI system during clinical procedures involving investigational drugs, with mandatory audit recording. Activation of an emergency stop transitions the task order to a failed state and requires a new safety check cycle before any subsequent procedure.

(k) \textit{Human oversight} means the requirement that qualified clinical personnel maintain supervisory authority over Physical AI systems during all phases of clinical investigation under an IND. Human oversight encompasses the ability to monitor, intervene, override, or terminate Physical AI system operations at any time. The human oversight quality management framework (v1.2.0) specifies risk-tiered controls with escalating requirements based on clinical consequence.

(l) \textit{Hash-chained audit trail} means an immutable, tamper-evident record of all Physical AI system actions using SHA-256 hash chains. Each audit record contains the action performed, actor identity, timestamp, resource affected, and a hash linking it to the previous record. This mechanism satisfies the electronic records requirements of 21 CFR Part 11 for IND-related activities.

(m) \textit{Model Context Protocol (MCP)} means an open protocol for connecting AI agents to external tools and data sources, implemented through five standardized servers: trialmcp-authz (authorization), trialmcp-fhir (clinical data), trialmcp-dicom (imaging), trialmcp-ledger (audit), and trialmcp-provenance (data lineage). The protocol defines 23 tools and operates under the governance of the Linux Foundation AI and Data Foundation.

(n) \textit{Deny-by-default authorization} means a security model where all Physical AI system access to clinical data associated with an IND is denied unless explicitly permitted by policy. Token-based sessions with defined scopes and time limits govern every data access request. This authorization model is enforced by the trialmcp-authz server.

(o) \textit{HIPAA Safe Harbor de-identification} means the removal or generalization of all 18 Protected Health Information (PHI) identifiers before clinical data is provided to Physical AI systems. HMAC-SHA256 with site-specific salts is used for pseudonymization of subject identifiers in IND-related records processed by Physical AI systems.

(p) \textit{Predetermined Change Control Plan (PCCP)} means a documented plan for managing modifications to Physical AI systems used as or within medical devices in connection with investigational drug administration, per FDA AI/ML guidance. The PCCP specifies the types of anticipated modifications, the methodology for implementing and validating changes, and the assessment of the impact on safety and effectiveness.

(q) \textit{Federated learning} means multi-site AI model training without sharing raw patient data, using differential privacy (epsilon and delta parameters) and secure aggregation. The federation framework supports FedAvg, FedProx, and SCAFFOLD algorithms with data harmonization for DICOM, FHIR, ICD-10 to SNOMED CT, and LOINC tumor marker standardization. Federated learning in IND-covered investigations must comply with the data protection requirements of this part.

(r) \textit{Pre-procedure safety matrix} means the set of verification checks that must pass before a Physical AI system may begin a clinical procedure involving an investigational drug. The matrix includes: valid authorization token, patient identity verification, clinical data availability (FHIR), imaging data access (DICOM), robot readiness and USL threshold verification, environmental checks, simulation validation confirmation, and digital twin synchronization status.

(s) \textit{Conformance level} means one of five cumulative levels defined in the national MCP-PAI standard for MCP server implementations: Core (base audit and authorization), Clinical Read (FHIR R4 clinical data access), Imaging (DICOM imaging with RECIST), Federated Site (multi-site federated learning with differential privacy), and Robot Procedure (full robot agent workflow with safety matrix and USL integration).

(t) \textit{Sim-to-real gap} means the difference between the behavior of a Physical AI system in simulation and its behavior in the physical world. For IND-covered investigations, the sim-to-real gap must be quantified using the validation benchmark suite (v1.0) and must demonstrate trajectory deviations below 2mm and force discrepancies below 0.5N for the specific clinical procedures to be performed.

(u) \textit{Cross-framework validation} means the verification that a Physical AI system's behavior is consistent across multiple simulation frameworks. Cross-framework validation requires testing in at least two recognized simulation frameworks with behavioral consistency metrics within the tolerances specified in the physical-ai-oncology-trials repository (v2.5.0).

(v) \textit{Vision-language-action (VLA) model} means a generative AI model that accepts visual and language inputs and produces robotic action outputs. VLA models used in Physical AI systems for IND-covered investigations must be validated for the specific clinical procedures to be performed, with performance metrics documented in the IND submission.

[52 FR 8831, Mar. 19, 1987, as amended at 64 FR 401, Jan. 5, 1999; 64 FR 56449, Oct. 20, 1999; 73 FR 22815, Apr. 28, 2008. Physical AI definitions added 17 March 2026.]

\subsection{\S~312.6 Labeling of an Investigational New Drug}

(a) The immediate package of an investigational new drug intended for human use shall bear a label with the statement ``Caution: New Drug---Limited by Federal (or United States) law to investigational use.''

(b) The label or labeling of an investigational new drug shall not bear any statement that is false or misleading in any particular and shall not represent that the investigational new drug is safe or effective for the purposes for which it is being investigated.

(c) The appropriate FDA Center Director, according to the procedures set forth in \S\S~201.26 or 610.68 of this chapter, may grant an exception or alternative to the provision in paragraph (a) of this section, to the extent that this provision is not explicitly required by statute, for specified lots, batches, or other units of a human drug product that is or will be included in the Strategic National Stockpile.

\subsubsection{Physical AI Adaptation of \S~312.6}

(d) When an investigational new drug is administered, dispensed, or monitored by a Physical AI system, the labeling of the investigational drug shall additionally include, either on the label or in accompanying documentation provided to the Physical AI system operator:

\begin{description}[leftmargin=0.8cm,style=sameline]
\item[(i)] A statement identifying the drug as an investigational new drug under an active IND;
\item[(ii)] The specific Physical AI system(s) authorized to administer, dispense, or monitor the drug, identified by system type, manufacturer, and model;
\item[(iii)] Any Physical AI system-specific administration parameters, including dosing rate constraints, positional accuracy requirements, and timing specifications that differ from manual administration; and
\item[(iv)] Machine-readable identifiers (barcode, RFID, or equivalent) sufficient to enable the Physical AI system to verify drug identity, lot number, and expiration date prior to administration.
\end{description}

(e) The labeling of a Physical AI system component that is integral to the administration of an investigational new drug (such as a robotic injector cartridge, automated infusion adapter, or needle-placement guide) shall bear the statement ``Caution: For use with investigational new drug under IND---Physical AI system component.''

[52 FR 8831, Mar. 19, 1987, as amended at 72 FR 73599, Dec. 28, 2007. Physical AI adaptation added 17 March 2026.]

\subsection{\S~312.7 Promotion of Investigational Drugs}

(a) \textit{Promotion of an investigational new drug.} A sponsor or investigator, or any person acting on behalf of a sponsor or investigator, shall not represent in a promotional context that an investigational new drug is safe or effective for the purposes for which it is under investigation or otherwise promote the drug. This provision is not intended to restrict the full exchange of scientific information concerning the drug, including dissemination of scientific findings in scientific or lay media. Rather, its intent is to restrict promotional claims of safety or effectiveness of the drug for a use for which it is under investigation and to preclude commercialization of the drug before it is approved for commercial distribution.

(b) \textit{Commercial distribution of an investigational new drug.} A sponsor or investigator shall not commercially distribute or test market an investigational new drug.

(c) \textit{Prolonging an investigation.} A sponsor shall not unduly prolong an investigation after finding that the results of the investigation appear to establish sufficient data to support a marketing application.

\subsubsection{Physical AI Adaptation of \S~312.7}

(d) \textit{Physical AI system promotion restrictions.} The promotion restrictions in paragraphs (a) through (c) of this section extend to claims regarding the performance, accuracy, or superiority of Physical AI systems used in connection with the administration of investigational new drugs. A sponsor or investigator shall not represent, in a promotional context, that a Physical AI system's involvement in drug administration enhances the safety or effectiveness of an investigational drug beyond what the clinical evidence supports.

(e) \textit{Simulation and digital twin data.} The dissemination of simulation results, digital twin predictions, or AI model performance metrics generated during an IND-covered investigation is permitted as part of the full exchange of scientific information, provided such dissemination does not constitute a promotional claim of drug safety or effectiveness. Simulation data presented in scientific communications shall be clearly labeled as computational predictions subject to validation.

[52 FR 8831, Mar. 19, 1987, as amended at 52 FR 19476, May 22, 1987; 67 FR 9585, Mar. 4, 2002; 74 FR 40899, Aug. 13, 2009. Physical AI adaptation added 17 March 2026.]

\subsection{\S~312.8 Charging for Investigational Drugs Under an IND}

(a) \textit{General criteria for charging.}

\begin{description}[leftmargin=0.5cm,style=sameline]
\item[(1)] A sponsor must meet the applicable requirements in paragraph (b) of this section for charging in a clinical trial or paragraph (c) of this section for charging for expanded access to an investigational drug for treatment use under subpart I of this part, except that sponsors need not fulfill the requirements in this section to charge for an approved drug obtained from another entity not affiliated with the sponsor for use as part of the clinical trial evaluation (e.g., in a clinical trial of a new use of the approved drug, for use of the approved drug as an active control).
\item[(2)] A sponsor must justify the amount to be charged in accordance with paragraph (d) of this section.
\item[(3)] A sponsor must obtain prior written authorization from FDA to charge for an investigational drug.
\item[(4)] FDA will withdraw authorization to charge if it determines that charging is interfering with the development of a drug for marketing approval or that the criteria for the authorization are no longer being met.
\end{description}

(b) \textit{Charging in a clinical trial.}

\begin{description}[leftmargin=0.5cm,style=sameline]
\item[(1)] \textit{Charging for a sponsor's drug.} A sponsor who wishes to charge for its investigational drug, including investigational use of its approved drug, must:
\begin{description}[leftmargin=0.8cm,style=sameline]
\item[(i)] Provide evidence that the drug has a potential clinical benefit that, if demonstrated in the clinical investigations, would provide a significant advantage over available products in the diagnosis, treatment, mitigation, or prevention of a disease or condition;
\item[(ii)] Demonstrate that the data to be obtained from the clinical trial would be essential to establishing that the drug is effective or safe for the purpose of obtaining initial approval of a drug, or would support a significant change in the labeling of an approved drug; and
\item[(iii)] Demonstrate that the clinical trial could not be conducted without charging because the cost of the drug is extraordinary to the sponsor.
\end{description}
\item[(2)] \textit{Duration of charging in a clinical trial.} Unless FDA specifies a shorter period, charging may continue for the length of the clinical trial.
\end{description}

(c) \textit{Charging for expanded access to investigational drug for treatment use.}

\begin{description}[leftmargin=0.5cm,style=sameline]
\item[(1)] A sponsor who wishes to charge for expanded access to an investigational drug for treatment use under subpart I of this part must provide reasonable assurance that charging will not interfere with developing the drug for marketing approval.
\item[(2)] For expanded access under \S~312.320 (treatment IND or treatment protocol), such assurance must include:
\begin{description}[leftmargin=0.8cm,style=sameline]
\item[(i)] Evidence of sufficient enrollment in any ongoing clinical trial(s) needed for marketing approval;
\item[(ii)] Evidence of adequate progress in the development of the drug for marketing approval; and
\item[(iii)] Information submitted under the general investigational plan (\S~312.23(a)(3)(iv)) specifying the drug development milestones the sponsor plans to meet in the next year.
\end{description}
\item[(3)] The authorization to charge is limited to the number of patients authorized to receive the drug under the treatment use, if there is a limitation.
\item[(4)] Unless FDA specifies a shorter period, charging for expanded access to an investigational drug for treatment use under subpart I of this part may continue for 1 year from the time of FDA authorization. A sponsor may request that FDA reauthorize charging for additional periods.
\end{description}

(d) \textit{Costs recoverable when charging for an investigational drug.}

\begin{description}[leftmargin=0.5cm,style=sameline]
\item[(1)] A sponsor may recover only the direct costs of making its investigational drug available.
\begin{description}[leftmargin=0.8cm,style=sameline]
\item[(i)] Direct costs are costs incurred by a sponsor that can be specifically and exclusively attributed to providing the drug for the investigational use for which FDA has authorized cost recovery. Direct costs include costs per unit to manufacture the drug or costs to acquire the drug from another manufacturing source, and direct costs to ship and handle the drug.
\item[(ii)] Indirect costs include costs incurred primarily to produce the drug for commercial sale and research and development, administrative, labor, or other costs that would be incurred even if the clinical trial or treatment use for which charging is authorized did not occur.
\end{description}
\item[(2)] For expanded access to an investigational drug for treatment use under \S\S~312.315 and 312.320, in addition to the direct costs described in paragraph (d)(1)(i) of this section, a sponsor may recover the costs of monitoring the expanded access IND or protocol, complying with IND reporting requirements, and other administrative costs directly associated with the expanded access IND.
\item[(3)] To support its calculation for cost recovery, a sponsor must provide supporting documentation to show that the calculation is consistent with the requirements of paragraphs (d)(1) and, if applicable, (d)(2) of this section. The documentation must be accompanied by a statement that an independent certified public accountant has reviewed and approved the calculations.
\end{description}

\subsubsection{Physical AI Adaptation of \S~312.8}

(e) \textit{Physical AI system costs.} When an investigational drug is administered through or monitored by a Physical AI system, the following cost considerations apply:

\begin{description}[leftmargin=0.8cm,style=sameline]
\item[(i)] The direct costs of Physical AI system deployment, including robotic hardware calibration, simulation validation, and digital twin configuration specific to the investigational drug, are not recoverable as drug costs under paragraph (d) of this section unless the Physical AI system is integral to the drug's mechanism of delivery and cannot be separated from the drug administration process.
\item[(ii)] Costs associated with Physical AI system safety validation, including pre-procedure safety matrix verification, USL assessment, and cross-framework simulation testing, are considered investigation infrastructure costs and are not recoverable as direct drug costs.
\item[(iii)] When a Physical AI system is classified as a combination product with the investigational drug, the sponsor may seek FDA guidance on the allocation of recoverable costs between the drug component and the Physical AI system component.
\end{description}

[74 FR 40899, Aug. 13, 2009. Physical AI adaptation added 17 March 2026.]

\subsection{\S~312.10 Waivers}

(a) A sponsor may request FDA to waive applicable requirement under this part. A waiver request may be submitted either in an IND or in an information amendment to an IND. In an emergency, a request may be made by telephone or other rapid communication means. A waiver request is required to contain at least one of the following:

\begin{description}[leftmargin=0.5cm,style=sameline]
\item[(1)] An explanation why the sponsor's compliance with the requirement is unnecessary or cannot be achieved;
\item[(2)] A description of an alternative submission or course of action that satisfies the purpose of the requirement; or
\item[(3)] Other information justifying a waiver.
\end{description}

(b) FDA may grant a waiver if it finds that the sponsor's noncompliance would not pose a significant and unreasonable risk to human subjects of the investigation and that one of the following is met:

\begin{description}[leftmargin=0.5cm,style=sameline]
\item[(1)] The sponsor's compliance with the requirement is unnecessary for the agency to evaluate the application, or compliance cannot be achieved;
\item[(2)] The sponsor's proposed alternative satisfies the requirement; or
\item[(3)] The applicant's submission otherwise justifies a waiver.
\end{description}

\subsubsection{Physical AI Adaptation of \S~312.10}

(c) \textit{Physical AI waiver provisions.} The waiver provisions of paragraphs (a) and (b) of this section apply to Physical AI system requirements established under this part. A sponsor seeking a waiver of Physical AI-specific requirements must, in addition to the requirements of paragraph (a):

\begin{description}[leftmargin=0.8cm,style=sameline]
\item[(i)] Demonstrate that the waiver would not compromise the safety of human subjects who interact with or are treated by Physical AI systems during the investigation;
\item[(ii)] Provide a risk assessment addressing the specific Physical AI system risks that the waived requirement was designed to mitigate; and
\item[(iii)] Describe alternative safety measures that will be implemented in lieu of the waived requirement.
\end{description}

(d) \textit{USL threshold waivers.} FDA may grant a waiver of the minimum USL threshold requirements specified in \S~312.1(e) if the sponsor demonstrates that the Physical AI system, despite not meeting the minimum USL score, has been validated through alternative means that provide equivalent assurance of safety and effectiveness for the specific clinical procedures to be performed. Such alternative validation must be documented and submitted with the waiver request.

(e) \textit{Simulation framework waivers.} FDA may grant a waiver of the cross-framework validation requirement if the sponsor demonstrates that the Physical AI system has been validated in a single simulation framework with sufficient fidelity and that cross-framework validation is not technically feasible for the specific system or procedure. The sponsor must provide evidence that the single-framework validation adequately characterizes the sim-to-real gap for the intended clinical use.

[52 FR 8831, Mar. 19, 1987, as amended at 52 FR 23031, June 17, 1987; 67 FR 9585, Mar. 4, 2002. Physical AI adaptation added 17 March 2026.]

\clearpage

\section{SUBPART B --- INVESTIGATIONAL NEW DRUG APPLICATION (IND)}

\subsection{\S~312.20 Requirement for an IND}

(a) A sponsor shall submit an IND to FDA if the sponsor intends to conduct a clinical investigation with an investigational new drug that is subject to \S~312.2(a).

(b) A sponsor shall not begin a clinical investigation subject to \S~312.2(a) until the investigation is subject to an IND which is in effect in accordance with \S~312.40.

(c) A sponsor shall submit a separate IND for any clinical investigation involving an exception from informed consent under \S~50.24 of this chapter. Such a clinical investigation is not permitted to proceed without the prior written authorization from FDA. FDA shall provide a written determination 30 days after FDA receives the IND or earlier.

\subsubsection{Physical AI Adaptation of \S~312.20}

(d) When a clinical investigation involves a Physical AI system that administers, dispenses, monitors, or otherwise participates in the investigational drug delivery pathway, the sponsor shall include in the IND a Physical AI System Description as specified in \S~312.23. The Physical AI System Description shall be submitted as part of the initial IND or, if the Physical AI system is introduced after the IND is in effect, as a protocol amendment under \S~312.30.

(e) A sponsor shall submit a separate protocol amendment for each distinct Physical AI system type (surgical robot, cobot, humanoid, therapeutic, diagnostic, assistive, or rehabilitative) to be used in the clinical investigation, unless the systems share substantially similar safety profiles and operational parameters, in which case they may be covered under a single protocol amendment with individual system specifications appended.

(f) When a Physical AI system is used in an investigation involving an exception from informed consent under \S~50.24 of this chapter, the IND must additionally document the Physical AI system's emergency deployment protocols, including rapid activation procedures, pre-loaded safety parameters, and the minimum human oversight requirements during emergency administration of the investigational drug.

[52 FR 8831, Mar. 19, 1987, as amended at 61 FR 51529, Oct. 2, 1996; 62 FR 32479, June 16, 1997. Physical AI adaptation added 17 March 2026.]

\subsection{\S~312.21 Phases of an Investigation}

An IND may be submitted for one or more phases of an investigation. The clinical investigation of a previously untested drug is generally divided into three phases. Although in general the phases are conducted sequentially, they may overlap. These three phases of an investigation are as follows:

(a) \textit{Phase 1.}

\begin{description}[leftmargin=0.5cm,style=sameline]
\item[(1)] Phase 1 includes the initial introduction of an investigational new drug into humans. Phase 1 studies are typically closely monitored and may be conducted in patients or normal volunteer subjects. These studies are designed to determine the metabolism and pharmacologic actions of the drug in humans, the side effects associated with increasing doses, and, if possible, to gain early evidence on effectiveness. During Phase 1, sufficient information about the drug's pharmacokinetics and pharmacological effects should be obtained to permit the design of well-controlled, scientifically valid, Phase 2 studies. The total number of subjects and patients included in Phase 1 studies varies with the drug, but is generally in the range of 20 to 80.
\item[(2)] Phase 1 studies also include studies of drug metabolism, structure-activity relationships, and mechanism of action in humans, as well as studies in which investigational drugs are used as research tools to explore biological phenomena or disease processes.
\end{description}

(b) \textit{Phase 2.} Phase 2 includes the controlled clinical studies conducted to evaluate the effectiveness of the drug for a particular indication or indications in patients with the disease or condition under study and to determine the common short-term side effects and risks associated with the drug. Phase 2 studies are typically well controlled, closely monitored, and conducted in a relatively small number of patients, usually involving no more than several hundred subjects.

(c) \textit{Phase 3.} Phase 3 studies are expanded controlled and uncontrolled trials. They are performed after preliminary evidence suggesting effectiveness of the drug has been obtained, and are intended to gather the additional information about effectiveness and safety that is needed to evaluate the overall benefit-risk relationship of the drug and to provide an adequate basis for physician labeling. Phase 3 studies usually include from several hundred to several thousand subjects.

\subsubsection{Physical AI Adaptation of \S~312.21}

(d) \textit{Phase 0 --- Simulation Validation.} For clinical investigations involving Physical AI systems, a simulation validation phase (Phase 0) shall precede human subject enrollment. Phase 0 consists of comprehensive testing of the Physical AI system in validated simulation environments using patient-derived digital twin models. Phase 0 is not a regulatory phase requiring separate IND submission but shall be documented in the IND as part of the pharmacology and toxicology information under \S~312.23(a)(8). Phase 0 validation shall include:

\begin{description}[leftmargin=0.8cm,style=sameline]
\item[(i)] Testing in at least two recognized simulation frameworks (NVIDIA Isaac Lab v2.3.1, MuJoCo v3.4.0, Gazebo Sim v10.0.0, or PyBullet v3.2.5) demonstrating cross-framework behavioral consistency with trajectory deviations below 2mm and force discrepancies below 0.5N;
\item[(ii)] Sim-to-real gap quantification using the validation benchmark suite with documented metrics for physics accuracy, contact dynamics, and task completion rates;
\item[(iii)] Digital twin-based procedure rehearsal using patient-specific anatomical models calibrated to clinical imaging data; and
\item[(iv)] Failure mode analysis across at least 1,000 simulated procedures documenting emergency stop activation rates, collision events, and force limit exceedances.
\end{description}

(e) \textit{Physical AI considerations by phase.}

\begin{description}[leftmargin=0.8cm,style=sameline]
\item[(i)] \textit{Phase 1 Physical AI requirements.} When Physical AI systems are used in Phase 1 investigations, the sponsor shall implement maximum human oversight with qualified operators maintaining direct manual override capability at all times. Physical AI systems used in Phase 1 shall operate in teleoperation or supervised autonomous mode only. Fully autonomous operation is not permitted in Phase 1.
\item[(ii)] \textit{Phase 2 Physical AI requirements.} Phase 2 investigations may permit graduated autonomy of Physical AI systems based on Phase 1 safety data. The protocol shall specify the autonomy level permitted and the conditions under which autonomy may be increased or decreased. USL scores shall be reassessed at the Phase 1 to Phase 2 transition using updated clinical performance data.
\item[(iii)] \textit{Phase 3 Physical AI requirements.} Phase 3 investigations may permit the full range of Physical AI autonomy levels supported by the accumulated safety data from Phases 1 and 2. Multi-site deployment of Physical AI systems in Phase 3 shall comply with the federated learning and cross-site coordination requirements of this part. Each site shall independently verify the Physical AI system's USL score and pre-procedure safety matrix before patient enrollment begins.
\end{description}

[52 FR 8831, Mar. 19, 1987. Physical AI adaptation added 17 March 2026.]

\subsection{\S~312.22 General Principles of the IND Submission}

(a) FDA's primary objectives in reviewing an IND are, in all phases of the investigation, to assure the safety and rights of subjects, and, in Phase 2 and 3, to help assure that the quality of the scientific evaluation of drugs is adequate to permit an evaluation of the drug's effectiveness and safety. Therefore, although FDA's review of Phase 1 submissions will focus on assessing the safety of Phase 1 investigations, FDA's review of Phases 2 and 3 submissions will also include an assessment of the scientific quality of the clinical investigations and the likelihood that the investigations will yield data capable of meeting statutory standards for marketing approval.

(b) The amount of information on a particular drug that must be submitted in an IND to assure the accomplishment of the objectives described in paragraph (a) of this section depends upon such factors as the novelty of the drug, the extent to which it has been studied previously, the known or suspected risks, and the developmental phase of the drug.

(c) The central focus of the initial IND submission should be on the general investigational plan and the protocols for specific human studies. Subsequent amendments to the IND that contain new or revised protocols should build logically on previous submissions and should be supported by additional information, including the results of animal toxicology studies or other human studies as appropriate. Annual reports to the IND should serve as the focus for reporting the status of studies being conducted under the IND and should update the general investigational plan for the coming year.

(d) The IND format set forth in \S~312.23 should be followed routinely by sponsors in the interest of fostering an efficient review of applications. Sponsors are expected to exercise considerable discretion, however, regarding the content of information submitted in each section, depending upon the kind of drug being studied and the nature of the available information. \S~312.23 outlines the information needed for a commercially sponsored IND for a new molecular entity. A sponsor-investigator who uses, as a research tool, an investigational new drug that is already subject to a manufacturer's IND or marketing application should follow the same general format, but ordinarily may, if authorized by the manufacturer, refer to the manufacturer's IND or marketing application in providing the technical information supporting the proposed clinical investigation. A sponsor-investigator who uses an investigational drug not subject to a manufacturer's IND or marketing application is ordinarily required to submit all technical information supporting the IND, unless such information may be referenced from the scientific literature.

\subsubsection{Physical AI Adaptation of \S~312.22}

(e) \textit{Physical AI submission principles.} In addition to the factors listed in paragraph (b) of this section, the amount of Physical AI system information required in an IND depends upon the complexity of the Physical AI system, the degree of autonomy involved, the criticality of the system's role in drug administration, and the maturity of the system as measured by its USL score. Systems with USL scores in the Advanced band (7.0 to 10.0) may require less extensive simulation validation documentation than systems in the Foundational band (1.0 to 4.9), provided the higher USL score is supported by documented clinical deployment history.

(f) \textit{Scalable documentation.} The Physical AI documentation requirements of \S~312.23 are scalable based on the role of the Physical AI system in the investigation:

\begin{description}[leftmargin=0.8cm,style=sameline]
\item[(i)] \textit{Primary role} (the Physical AI system directly administers or dispenses the investigational drug): Full Physical AI System Description, simulation validation data, digital twin specifications, and cybersecurity assessment required.
\item[(ii)] \textit{Secondary role} (the Physical AI system monitors or assists with drug administration under direct human control): Abbreviated Physical AI System Description focusing on safety parameters, human-robot interaction protocols, and monitoring accuracy.
\item[(iii)] \textit{Ancillary role} (the Physical AI system performs supportive tasks such as sample handling, logistics, or patient positioning that do not directly involve drug contact): Summary Physical AI System Description with safety parameters and human oversight protocols.
\end{description}

[52 FR 8831, Mar. 19, 1987. Physical AI adaptation added 17 March 2026.]

\subsection{\S~312.23 IND Content and Format}

(a) A sponsor who intends to conduct a clinical investigation subject to this part shall submit an ``Investigational New Drug Application'' (IND) including, in the following order:

\begin{description}[leftmargin=0.5cm,style=sameline]
\item[(1)] \textit{Cover sheet (Form FDA-1571).} A cover sheet for the application containing the following:
\begin{description}[leftmargin=0.8cm,style=sameline]
\item[(i)] The name, address, and telephone number of the sponsor, the date of the application, and the name of the investigational new drug.
\item[(ii)] Identification of the phase or phases of the clinical investigation to be conducted.
\item[(iii)] A commitment not to begin clinical investigations until an IND covering the investigations is in effect.
\item[(iv)] A commitment that an Institutional Review Board (IRB) that complies with the requirements set forth in part 56 will be responsible for the initial and continuing review and approval of each of the studies in the proposed clinical investigation and that the investigator will report to the IRB proposed changes in the research activity in accordance with the requirements of part 56.
\item[(v)] A commitment to conduct the investigation in accordance with all other applicable regulatory requirements.
\item[(vi)] The name and title of the person responsible for monitoring the conduct and progress of the clinical investigations.
\item[(vii)] The name(s) and title(s) of the person(s) responsible under \S~312.32 for review and evaluation of information relevant to the safety of the drug.
\item[(viii)] If a sponsor has transferred any obligations for the conduct of any clinical study to a contract research organization, a statement containing the name and address of the contract research organization, identification of the clinical study, and a listing of the obligations transferred.
\item[(ix)] The signature of the sponsor or the sponsor's authorized representative.
\end{description}

\item[(2)] \textit{A table of contents.}

\item[(3)] \textit{Introductory statement and general investigational plan.}
\begin{description}[leftmargin=0.8cm,style=sameline]
\item[(i)] A brief introductory statement giving the name of the drug and all active ingredients, the drug's pharmacological class, the structural formula of the drug (if known), the formulation of the dosage form(s) to be used, the route of administration, and the broad objectives and planned duration of the proposed clinical investigation(s).
\item[(ii)] A brief summary of previous human experience with the drug, with reference to other IND's if pertinent, and to investigational or marketing experience in other countries that may be relevant to the safety of the proposed clinical investigation(s).
\item[(iii)] If the drug has been withdrawn from investigation or marketing in any country for any reason related to safety or effectiveness, identification of the country(ies) where the drug was withdrawn and the reasons for the withdrawal.
\item[(iv)] A brief description of the overall plan for investigating the drug product for the following year, including: (\textit{a}) the rationale for the drug or the research study; (\textit{b}) the indication(s) to be studied; (\textit{c}) the general approach to be followed in evaluating the drug; (\textit{d}) the kinds of clinical trials to be conducted in the first year following the submission; (\textit{e}) the estimated number of patients to be given the drug in those studies; and (\textit{f}) any risks of particular severity or seriousness anticipated on the basis of the toxicological data in animals or prior studies in humans with the drug or related drugs.
\end{description}

\item[(4)] [Reserved]

\item[(5)] \textit{Investigator's brochure.} If required under \S~312.55, a copy of the investigator's brochure, containing the following information:
\begin{description}[leftmargin=0.8cm,style=sameline]
\item[(i)] A brief description of the drug substance and the formulation, including the structural formula, if known.
\item[(ii)] A summary of the pharmacological and toxicological effects of the drug in animals and, to the extent known, in humans.
\item[(iii)] A summary of the pharmacokinetics and biological disposition of the drug in animals and, if known, in humans.
\item[(iv)] A summary of information relating to safety and effectiveness in humans obtained from prior clinical studies.
\item[(v)] A description of possible risks and side effects to be anticipated on the basis of prior experience with the drug under investigation or with related drugs, and of precautions or special monitoring to be done as part of the investigational use of the drug.
\end{description}

\item[(6)] \textit{Protocols.}
\begin{description}[leftmargin=0.8cm,style=sameline]
\item[(i)] A protocol for each planned study. In general, protocols for Phase 1 studies may be less detailed and more flexible than protocols for Phase 2 and 3 studies. Phase 1 protocols should be directed primarily at providing an outline of the investigation---an estimate of the number of patients to be involved, a description of safety exclusions, and a description of the dosing plan including duration, dose, or method to be used in determining dose---and should specify in detail only those elements of the study that are critical to safety, such as necessary monitoring of vital signs and blood chemistries. Modifications of the experimental design of Phase 1 studies that do not affect critical safety assessments are required to be reported to FDA only in the annual report.
\item[(ii)] In Phases 2 and 3, detailed protocols describing all aspects of the study should be submitted.
\item[(iii)] A protocol is required to contain the following: (\textit{a}) a statement of the objectives and purpose of the study; (\textit{b}) the name, address, and qualifications of each investigator and subinvestigator, the research facilities, and the reviewing IRB; (\textit{c}) the criteria for patient selection and exclusion and an estimate of the number of patients; (\textit{d}) a description of the design of the study, including the control group and methods to minimize bias; (\textit{e}) the method for determining the dose(s), the planned maximum dosage, and the duration of individual patient exposure; (\textit{f}) a description of the observations and measurements to fulfill the study objectives; and (\textit{g}) a description of clinical procedures, laboratory tests, or other measures to monitor the effects of the drug and to minimize risk.
\end{description}

\item[(7)] \textit{Chemistry, manufacturing, and control information.}
\begin{description}[leftmargin=0.8cm,style=sameline]
\item[(i)] As appropriate for the particular investigations covered by the IND, a section describing the composition, manufacture, and control of the drug substance and the drug product. The amount of information needed will vary with the phase of the investigation, the proposed duration, the dosage form, and the amount of information otherwise available.
\item[(ii)] The amount of information to be submitted depends upon the scope of the proposed clinical investigation.
\item[(iii)] As drug development proceeds and as the scale of production is changed, the sponsor should submit information amendments to supplement the initial information.
\item[(iv)] Based on the phase(s) to be studied, the submission is required to contain: (\textit{a}) \textit{Drug substance}: a description including physical, chemical, or biological characteristics; manufacturer name and address; general method of preparation; acceptable limits and analytical methods; and stability information. (\textit{b}) \textit{Drug product}: a list of all components, quantitative composition, manufacturer name and address, manufacturing and packaging procedure, acceptable limits and analytical methods, and stability information. (\textit{c}) A brief general description of the composition, manufacture, and control of any placebo. (\textit{d}) \textit{Labeling}: a copy of all labels and labeling to be provided to each investigator. (\textit{e}) \textit{Environmental analysis requirements}: a claim for categorical exclusion under \S~25.30 or 25.31 or an environmental assessment under \S~25.40.
\end{description}

\item[(8)] \textit{Pharmacology and toxicology information.} Adequate information about pharmacological and toxicological studies of the drug involving laboratory animals or in vitro, on the basis of which the sponsor has concluded that it is reasonably safe to conduct the proposed clinical investigations. Such information is required to include the identification and qualifications of the individuals who evaluated the results and a statement of where the investigations were conducted and where the records are available for inspection.
\begin{description}[leftmargin=0.8cm,style=sameline]
\item[(i)] \textit{Pharmacology and drug disposition.} A section describing the pharmacological effects and mechanism(s) of action of the drug in animals, and information on the absorption, distribution, metabolism, and excretion of the drug, if known.
\item[(ii)] \textit{Toxicology.} (\textit{a}) An integrated summary of the toxicological effects of the drug in animals and in vitro, including acute, subacute, and chronic toxicity tests; tests of effects on reproduction and the developing fetus; and any special toxicity tests. (\textit{b}) For each toxicology study primarily supporting the safety of the proposed clinical investigation, a full tabulation of data suitable for detailed review.
\item[(iii)] For each nonclinical laboratory study subject to good laboratory practice regulations under part 58, a statement of compliance or a brief statement of the reason for noncompliance.
\end{description}

\item[(9)] \textit{Previous human experience with the investigational drug.} A summary of previous human experience known to the applicant, if any, with the investigational drug, including:
\begin{description}[leftmargin=0.8cm,style=sameline]
\item[(i)] Detailed information about previous investigations or marketing experience relevant to the safety of the proposed investigation.
\item[(ii)] If the drug is a combination of drugs previously investigated or marketed, the required information for each active drug component.
\item[(iii)] If the drug has been marketed outside the United States, a list of the countries in which it has been marketed and withdrawn from marketing for safety or effectiveness reasons.
\end{description}

\item[(10)] \textit{Additional information.}
\begin{description}[leftmargin=0.8cm,style=sameline]
\item[(i)] \textit{Drug dependence and abuse potential.} If the drug is a psychotropic substance or otherwise has abuse potential, relevant clinical studies and experience and studies in test animals.
\item[(ii)] \textit{Radioactive drugs.} If the drug is a radioactive drug, sufficient data for dosimetry calculations.
\item[(iii)] \textit{Pediatric studies.} Plans for assessing pediatric safety and effectiveness.
\item[(iv)] \textit{Other information.} A brief statement of any other information that would aid evaluation of the proposed clinical investigations.
\end{description}

\item[(11)] \textit{Relevant information.} If requested by FDA, any other relevant information needed for review of the application.
\end{description}

(b) \textit{Information previously submitted.} The sponsor ordinarily is not required to resubmit information previously submitted, but may incorporate the information by reference. A reference to information submitted previously must identify the file by name, reference number, volume, and page number.

(c) \textit{Material in a foreign language.} The sponsor shall submit an accurate and complete English translation of each part of the IND that is not in English.

(d) \textit{Number of copies.} The sponsor shall submit an original and two copies of all submissions to the IND file.

(e) \textit{Numbering of IND submissions.} Each submission relating to an IND is required to be numbered serially using a single, three-digit serial number. The initial IND is required to be numbered 000.

(f) \textit{Identification of exception from informed consent.} If the investigation involves an exception from informed consent under \S~50.24 of this chapter, the sponsor shall prominently identify on the cover sheet that the investigation is subject to the requirements in \S~50.24 of this chapter.

\subsubsection{Physical AI Adaptation of \S~312.23}

(g) \textit{Physical AI System Description.} When the clinical investigation involves a Physical AI system, the IND shall include, as a separate section following the information required by paragraph (a)(8), a Physical AI System Description containing the following:

\begin{description}[leftmargin=0.8cm,style=sameline]
\item[(i)] \textit{System identification and classification.} The manufacturer, model designation, software version, and robot category (surgical, cobot, humanoid, therapeutic, diagnostic, assistive, or rehabilitative) of each Physical AI system to be used. The current USL score for each system, computed across the four dimensions (simulation switching, AI integration, cross-robot sharing, clinical trial collaboration), with supporting documentation for the score calculation.

\item[(ii)] \textit{Role in drug administration.} A detailed description of the Physical AI system's role in the investigational drug pathway, classified as primary (direct drug administration or dispensing), secondary (monitoring or assisted administration under human control), or ancillary (supportive tasks not involving direct drug contact). The description shall include the specific clinical procedures to be performed, the degree of autonomy permitted, and the human oversight requirements for each procedure.

\item[(iii)] \textit{Hardware specifications.} Technical specifications of the robotic platform, including degrees of freedom, payload capacity, positional accuracy, force sensing resolution, workspace dimensions, and end-effector specifications relevant to drug administration. For surgical robots, the specifications shall include instrument tip accuracy, tremor compensation parameters, and tissue interaction force limits.

\item[(iv)] \textit{Software and AI/ML specifications.} A description of all software components, including the operating system, middleware (ROS 2 Jazzy or Kilted where applicable), control algorithms, and embedded AI/ML models. For each AI/ML model, the description shall include the model architecture, training data characteristics (including demographic diversity and oncology-specific representation), validation methodology, performance metrics, and the Predetermined Change Control Plan (PCCP) if the model is expected to be updated during the investigation.

\item[(v)] \textit{Simulation validation data.} The results of Phase 0 simulation validation as described in \S~312.21(d), including cross-framework behavioral consistency metrics, sim-to-real gap quantification, failure mode analysis results, and digital twin procedure rehearsal outcomes. The validation data shall reference the specific simulation framework versions used and the validation benchmark suite (v1.0) metrics achieved.

\item[(vi)] \textit{Digital twin specifications.} If digital twins are used for treatment planning or intraoperative guidance, a description of the digital twin modeling methodology, calibration data sources, update frequency, validation metrics, and the mathematical frameworks employed (reaction-diffusion equations, pharmacokinetic/pharmacodynamic models, finite element analysis). The description shall include the digital twin's role in the drug administration decision pathway.

\item[(vii)] \textit{Safety systems.} A comprehensive description of the Physical AI system's safety architecture, including emergency stop mechanisms (hardware and software), force and torque limits, collision detection and avoidance systems, workspace boundary enforcement, and fail-safe behaviors. The description shall include the pre-procedure safety matrix requirements and the automated verification checks performed before each clinical procedure.

\item[(viii)] \textit{Cybersecurity assessment.} An assessment of the Physical AI system's cybersecurity posture, including network architecture, authentication and authorization mechanisms (deny-by-default per the national MCP-PAI standard), encryption protocols, vulnerability management procedures, and incident response plans. The assessment shall address threats specific to robotic systems in clinical settings, including unauthorized command injection, sensor data manipulation, and communication interception.

\item[(ix)] \textit{MCP integration.} If the Physical AI system interfaces with clinical data through the Model Context Protocol, a description of the MCP server topology, the tools and resources accessed, the conformance level achieved (Core, Clinical Read, Imaging, Federated Site, or Robot Procedure), and the hash-chained audit trail implementation. The description shall include the robot capability profile (JSON schema) and the task order lifecycle management approach.

\item[(x)] \textit{Human-robot interaction protocols.} A description of the interfaces between the Physical AI system and clinical personnel, including operator control stations, status displays, alarm systems, communication channels, and handoff procedures. The description shall specify the training requirements for personnel who operate or supervise the Physical AI system during the investigation.

\item[(xi)] \textit{Maintenance and calibration.} The scheduled maintenance intervals, calibration procedures, and performance verification protocols for the Physical AI system. The description shall include the criteria for taking a system out of service and the re-qualification procedures required before returning the system to clinical use.
\end{description}

(h) \textit{Investigator's brochure Physical AI supplement.} When Physical AI systems are used in the investigation, the investigator's brochure required under paragraph (a)(5) shall include a Physical AI supplement containing:

\begin{description}[leftmargin=0.8cm,style=sameline]
\item[(i)] A plain-language summary of the Physical AI system's capabilities and limitations for investigator reference;
\item[(ii)] Known and anticipated risks specific to Physical AI system interactions, including robotic malfunction modes, AI prediction error rates, and human-robot interaction hazards;
\item[(iii)] Emergency procedures for Physical AI system failures, including manual override instructions and patient stabilization protocols; and
\item[(iv)] The minimum training requirements and competency assessments for investigators and subinvestigators who will operate or supervise the Physical AI system.
\end{description}

(i) \textit{Protocol Physical AI section.} When Physical AI systems are used in the investigation, each protocol submitted under paragraph (a)(6) shall include a Physical AI section containing:

\begin{description}[leftmargin=0.8cm,style=sameline]
\item[(i)] The specific Physical AI system(s) to be used in the study, referenced by the system identification in the Physical AI System Description;
\item[(ii)] The autonomy level permitted for each procedure type and the conditions under which autonomy may be increased, decreased, or revoked;
\item[(iii)] The Physical AI-specific inclusion and exclusion criteria for subjects (e.g., anatomical constraints, implant compatibility, claustrophobia assessment for enclosed robotic systems);
\item[(iv)] The Physical AI-specific observations and measurements, including system telemetry data, force profiles, trajectory accuracy, and procedure timing;
\item[(v)] The Physical AI-specific stopping rules, including system performance degradation thresholds, consecutive failure limits, and sim-to-real divergence criteria; and
\item[(vi)] The plan for federated learning coordination if the study involves multi-site deployment of Physical AI systems, including data harmonization protocols, differential privacy parameters, and aggregation algorithms.
\end{description}

[52 FR 8831, Mar. 19, 1987, as amended at 52 FR 23031, June 17, 1987; 53 FR 1918, Jan. 25, 1988; 61 FR 51529, Oct. 2, 1996; 62 FR 40599, July 29, 1997; 63 FR 66669, Dec. 2, 1998; 65 FR 56479, Sept. 19, 2000; 67 FR 9585, Mar. 4, 2002. Physical AI adaptation added 17 March 2026.]

\subsection{\S~312.30 Protocol Amendments}

Once an IND is in effect, a sponsor shall amend it as needed to ensure that the clinical investigations are conducted according to protocols included in the application. This section sets forth the provisions under which new protocols may be submitted and changes in previously submitted protocols may be made. Whenever a sponsor intends to conduct a clinical investigation with an exception from informed consent for emergency research as set forth in \S~50.24 of this chapter, the sponsor shall submit a separate IND for such investigation.

(a) \textit{New protocol.} Whenever a sponsor intends to conduct a study that is not covered by a protocol already contained in the IND, the sponsor shall submit to FDA a protocol amendment containing the protocol for the study. Such study may begin provided two conditions are met:

\begin{description}[leftmargin=0.5cm,style=sameline]
\item[(1)] The sponsor has submitted the protocol to FDA for its review; and
\item[(2)] the protocol has been approved by the Institutional Review Board (IRB) with responsibility for review and approval of the study in accordance with the requirements of part 56. The sponsor may comply with these two conditions in either order.
\end{description}

(b) \textit{Changes in a protocol.}

\begin{description}[leftmargin=0.5cm,style=sameline]
\item[(1)] A sponsor shall submit a protocol amendment describing any change in a Phase 1 protocol that significantly affects the safety of subjects or any change in a Phase 2 or 3 protocol that significantly affects the safety of subjects, the scope of the investigation, or the scientific quality of the study. Examples of changes requiring an amendment under this paragraph include:
\begin{description}[leftmargin=0.8cm,style=sameline]
\item[(i)] Any increase in drug dosage or duration of exposure of individual subjects to the drug beyond that in the current protocol, or any significant increase in the number of subjects under study.
\item[(ii)] Any significant change in the design of a protocol (such as the addition or dropping of a control group).
\item[(iii)] The addition of a new test or procedure that is intended to improve monitoring for, or reduce the risk of, a side effect or adverse event; or the dropping of a test intended to monitor safety.
\end{description}
\item[(2)] A protocol change under paragraph (b)(1) of this section may be made provided the sponsor has submitted the change to FDA and the change has been approved by the IRB. Notwithstanding the foregoing, a protocol change intended to eliminate an apparent immediate hazard to subjects may be implemented immediately provided FDA is subsequently notified by protocol amendment and the reviewing IRB is notified in accordance with \S~56.104(c).
\end{description}

(c) \textit{New investigator.} A sponsor shall submit a protocol amendment when a new investigator is added to carry out a previously submitted protocol, except that a protocol amendment is not required when a licensed practitioner is added in the case of a treatment protocol under \S~312.315 or \S~312.320. The sponsor shall notify FDA of the new investigator within 30 days.

(d) \textit{Content and format.} A protocol amendment is required to be prominently identified as such and to contain the following:

\begin{description}[leftmargin=0.5cm,style=sameline]
\item[(1)] In the case of a new protocol, a copy of the new protocol and a brief description of the most clinically significant differences between it and previous protocols. In the case of a change in protocol, a brief description of the change. In the case of a new investigator, the investigator's name, qualifications, reference to the previously submitted protocol, and all additional required information.
\item[(2)] Reference to specific technical information in the IND or in a concurrently submitted information amendment that the sponsor relies on to support any clinically significant change.
\item[(3)] If the sponsor desires FDA to comment on the submission, a request for such comment and the specific questions FDA's response should address.
\end{description}

(e) \textit{When submitted.} A sponsor shall submit a protocol amendment for a new protocol or a change in protocol before its implementation. Protocol amendments to add a new investigator or to provide additional information about investigators may be grouped and submitted at 30-day intervals.

\subsubsection{Physical AI Adaptation of \S~312.30}

(f) \textit{Physical AI protocol amendments.} In addition to the changes described in paragraph (b)(1) of this section, the following Physical AI system changes require a protocol amendment:

\begin{description}[leftmargin=0.8cm,style=sameline]
\item[(i)] Introduction of a new Physical AI system type or model not previously described in the Physical AI System Description under \S~312.23(g);
\item[(ii)] Any change in the autonomy level permitted for a Physical AI system during clinical procedures, including changes from teleoperation to supervised autonomy or from supervised autonomy to full autonomy;
\item[(iii)] Any significant update to an AI/ML model embedded in a Physical AI system that is outside the scope of the Predetermined Change Control Plan (PCCP), including changes to model architecture, training data composition, or performance thresholds;
\item[(iv)] Any change in the simulation framework or validation methodology used for ongoing Physical AI system verification;
\item[(v)] Any change in the human oversight requirements for Physical AI system operation, including changes to operator-to-system ratios, monitoring protocols, or intervention thresholds;
\item[(vi)] Any change in the digital twin modeling methodology, calibration data sources, or validation criteria that affects treatment planning or intraoperative guidance decisions; and
\item[(vii)] Any change in the MCP server topology, conformance level, or security architecture that affects the Physical AI system's access to clinical data.
\end{description}

(g) \textit{AI/ML model updates within PCCP scope.} Changes to AI/ML models embedded in Physical AI systems that fall within the scope of a previously submitted PCCP do not require a protocol amendment, provided the sponsor maintains documentation of all model updates and includes a summary of PCCP-covered changes in the annual report under \S~312.33. The sponsor shall notify FDA within 15 calendar days if a PCCP-covered update results in a performance metric falling below the pre-specified acceptance criteria.

(h) \textit{Emergency Physical AI modifications.} A Physical AI system modification intended to eliminate an apparent immediate safety hazard may be implemented immediately, consistent with paragraph (b)(2) of this section. The sponsor shall submit a protocol amendment describing the emergency modification within 15 working days and shall include an assessment of the modification's impact on system safety, USL score, and ongoing simulation validation.

[52 FR 8831, Mar. 19, 1987, as amended at 52 FR 23031, June 17, 1987; 53 FR 1918, Jan. 25, 1988; 61 FR 51530, Oct. 2, 1996; 67 FR 9585, Mar. 4, 2002; 74 FR 40942, Aug. 13, 2009. Physical AI adaptation added 17 March 2026.]

\subsection{\S~312.31 Information Amendments}

(a) \textit{Requirement for information amendment.} A sponsor shall report in an information amendment essential information on the IND that is not within the scope of a protocol amendment, IND safety reports, or annual report. Examples of information requiring an information amendment include:

\begin{description}[leftmargin=0.5cm,style=sameline]
\item[(1)] New toxicology, chemistry, or other technical information; or
\item[(2)] A report regarding the discontinuance of a clinical investigation.
\end{description}

(b) \textit{Content and format of an information amendment.} An information amendment is required to bear prominent identification of its contents and to contain: (1) a statement of the nature and purpose of the amendment; (2) an organized submission of the data in a format appropriate for scientific review; and (3) if the sponsor desires FDA to comment, a request for such comment.

(c) \textit{When submitted.} Information amendments to the IND should be submitted as necessary but, to the extent feasible, not more than every 30 days.

\subsubsection{Physical AI Adaptation of \S~312.31}

(d) \textit{Physical AI information amendments.} In addition to the examples in paragraph (a), the following Physical AI-related information requires an information amendment:

\begin{description}[leftmargin=0.8cm,style=sameline]
\item[(i)] New simulation validation data, including results from updated simulation framework versions or new cross-framework validation studies;
\item[(ii)] Updated USL scores for Physical AI systems used in the investigation, whether resulting from new assessment data or changes to the USL scoring methodology;
\item[(iii)] New cybersecurity vulnerability information affecting Physical AI systems used in the investigation, including vendor security advisories and patch deployment status;
\item[(iv)] Changes to the national MCP-PAI standard or the Physical AI system's MCP conformance level that affect the system's interaction with clinical data;
\item[(v)] New published literature or regulatory actions in other jurisdictions related to the safety or performance of Physical AI systems of the same type used in the investigation; and
\item[(vi)] Updates to the digital twin modeling framework, including new calibration data sources, algorithm improvements, or validation results that are not significant enough to require a protocol amendment.
\end{description}

[52 FR 8831, Mar. 19, 1987, as amended at 52 FR 23031, June 17, 1987; 53 FR 1918, Jan. 25, 1988; 67 FR 9585, Mar. 4, 2002. Physical AI adaptation added 17 March 2026.]

\subsection{\S~312.32 IND Safety Reporting}

(a) \textit{Definitions.} The following definitions of terms apply to this section:

\textit{Adverse event} means any untoward medical occurrence associated with the use of a drug in humans, whether or not considered drug related.

\textit{Life-threatening adverse event or life-threatening suspected adverse reaction.} An adverse event or suspected adverse reaction is considered ``life-threatening'' if, in the view of either the investigator or sponsor, its occurrence places the patient or subject at immediate risk of death. It does not include an adverse event or suspected adverse reaction that, had it occurred in a more severe form, might have caused death.

\textit{Serious adverse event or serious suspected adverse reaction.} An adverse event or suspected adverse reaction is considered ``serious'' if, in the view of either the investigator or sponsor, it results in any of the following outcomes: death, a life-threatening adverse event, inpatient hospitalization or prolongation of existing hospitalization, a persistent or significant incapacity or substantial disruption of the ability to conduct normal life functions, or a congenital anomaly/birth defect. Important medical events that may not result in death, be life-threatening, or require hospitalization may be considered serious when, based upon appropriate medical judgment, they may jeopardize the patient or subject and may require medical or surgical intervention to prevent one of the listed outcomes.

\textit{Suspected adverse reaction} means any adverse event for which there is a reasonable possibility that the drug caused the adverse event. For the purposes of IND safety reporting, ``reasonable possibility'' means there is evidence to suggest a causal relationship between the drug and the adverse event.

\textit{Unexpected adverse event or unexpected suspected adverse reaction.} An adverse event or suspected adverse reaction is considered ``unexpected'' if it is not listed in the investigator brochure or is not listed at the specificity or severity that has been observed; or, if an investigator brochure is not required or available, is not consistent with the risk information described in the general investigational plan or elsewhere in the current application, as amended.

(b) \textit{Review of safety information.} The sponsor must promptly review all information relevant to the safety of the drug obtained or otherwise received by the sponsor from foreign or domestic sources, including information derived from any clinical or epidemiological investigations, animal or in vitro studies, reports in the scientific literature, and unpublished scientific papers, as well as reports from foreign regulatory authorities and reports of foreign commercial marketing experience for drugs that are not marketed in the United States.

(c) \textit{IND safety reports.}

\begin{description}[leftmargin=0.5cm,style=sameline]
\item[(1)] The sponsor must notify FDA and all participating investigators in an IND safety report of potential serious risks, from clinical trials or any other source, as soon as possible, but in no case later than 15 calendar days after the sponsor determines that the information qualifies for reporting under paragraph (c)(1)(i), (c)(1)(ii), (c)(1)(iii), or (c)(1)(iv) of this section. In each IND safety report, the sponsor must identify all IND safety reports previously submitted concerning a similar suspected adverse reaction, and must analyze the significance of the suspected adverse reaction in light of previous reports or any other relevant information.

\begin{description}[leftmargin=0.8cm,style=sameline]
\item[(i)] \textit{Serious and unexpected suspected adverse reaction.} The sponsor must report any suspected adverse reaction that is both serious and unexpected, where there is evidence to suggest a causal relationship between the drug and the adverse event.
\item[(ii)] \textit{Findings from other studies.} The sponsor must report any findings from epidemiological studies, pooled analysis, or clinical studies that suggest a significant risk in humans exposed to the drug.
\item[(iii)] \textit{Findings from animal or in vitro testing.} The sponsor must report any findings from animal or in vitro testing that suggest a significant risk in humans exposed to the drug, such as reports of mutagenicity, teratogenicity, or carcinogenicity.
\item[(iv)] \textit{Increased rate of occurrence of serious suspected adverse reactions.} The sponsor must report any clinically important increase in the rate of a serious suspected adverse reaction over that listed in the protocol or investigator brochure.
\item[(v)] \textit{Submission of IND safety reports.} The sponsor must submit each IND safety report in a narrative format or on FDA Form 3500A or in an electronic format that FDA can process, review, and archive.
\end{description}

\item[(2)] \textit{Unexpected fatal or life-threatening suspected adverse reaction reports.} The sponsor must also notify FDA of any unexpected fatal or life-threatening suspected adverse reaction as soon as possible but in no case later than 7 calendar days after the sponsor's initial receipt of the information.

\item[(3)] \textit{Reporting format or frequency.} FDA may require a sponsor to submit IND safety reports in a format or at a frequency different than that required under this paragraph.

\item[(4)] \textit{Investigations of marketed drugs.} A sponsor of a clinical study of a drug marketed or approved in the United States that is conducted under an IND is required to submit IND safety reports for suspected adverse reactions observed in the clinical study.

\item[(5)] \textit{Reporting study endpoints.} Study endpoints (e.g., mortality or major morbidity) must be reported to FDA as described in the protocol and ordinarily would not be reported under paragraph (c) of this section, unless a serious and unexpected adverse event occurs for which there is evidence suggesting a causal relationship between the drug and the event.
\end{description}

(d) \textit{Followup.}

\begin{description}[leftmargin=0.5cm,style=sameline]
\item[(1)] The sponsor must promptly investigate all safety information it receives.
\item[(2)] Relevant followup information to an IND safety report must be submitted as soon as the information is available and must be identified as a ``Followup IND Safety Report.''
\item[(3)] If the results of a sponsor's investigation show that an adverse event not initially determined to be reportable under paragraph (c) of this section is so reportable, the sponsor must report it as soon as possible, but in no case later than 15 calendar days after the determination is made.
\end{description}

(e) \textit{Disclaimer.} A safety report or other information submitted by a sponsor under this part does not necessarily reflect a conclusion by the sponsor or FDA that the report or information constitutes an admission that the drug caused or contributed to an adverse event.

\subsubsection{Physical AI Adaptation of \S~312.32}

(f) \textit{Physical AI safety event definitions.} In addition to the definitions in paragraph (a), the following definitions apply to safety reporting for investigations involving Physical AI systems:

\textit{Physical AI adverse event} means any untoward occurrence associated with the operation of a Physical AI system during a clinical investigation involving an investigational drug, whether or not considered to be caused by the Physical AI system. Physical AI adverse events include robotic malfunction, AI prediction error, sensor failure, communication loss, unauthorized system access, digital twin divergence, and any event requiring activation of the emergency stop.

\textit{Serious Physical AI adverse event} means a Physical AI adverse event that results in, or has the potential to result in, any of the outcomes listed in the definition of serious adverse event in paragraph (a), or that results in unplanned interruption of an investigational drug administration procedure requiring clinical intervention, uncontrolled release of the investigational drug, or compromise of the Physical AI system's sterile field during a surgical procedure.

(g) \textit{Physical AI safety reporting requirements.} The sponsor must include Physical AI adverse events in IND safety reports as follows:

\begin{description}[leftmargin=0.8cm,style=sameline]
\item[(i)] \textit{Serious Physical AI adverse events.} Any serious Physical AI adverse event must be reported under the same timelines and procedures as paragraph (c)(1) of this section. The report must include the Physical AI system identification, the nature of the malfunction or error, the system telemetry data for the period surrounding the event, the emergency stop activation status, and the patient outcome.
\item[(ii)] \textit{Physical AI system-drug interaction events.} Any event in which a Physical AI system error directly or indirectly results in incorrect drug administration (wrong dose, wrong route, wrong timing, wrong patient) must be reported as an IND safety report regardless of whether the event meets the criteria for a serious adverse event.
\item[(iii)] \textit{Cybersecurity incidents.} Any cybersecurity incident affecting a Physical AI system during a clinical investigation, including unauthorized access attempts, data integrity violations, and communication interceptions, must be reported in an IND safety report within 15 calendar days of detection. If the incident results in or has the potential to result in patient harm, the 7-calendar-day reporting timeline for unexpected fatal or life-threatening events applies.
\item[(iv)] \textit{AI model performance degradation.} The sponsor must report any sustained degradation in AI/ML model performance below the pre-specified acceptance criteria documented in the PCCP or the Physical AI System Description. Sustained degradation is defined as performance below acceptance criteria for three or more consecutive procedures or for a cumulative period of 24 hours, whichever occurs first.
\item[(v)] \textit{Sim-to-real divergence.} The sponsor must report any instance in which the observed behavior of a Physical AI system during a clinical procedure diverges from the simulated behavior by more than twice the validated sim-to-real gap tolerance (i.e., trajectory deviation greater than 4mm or force discrepancy greater than 1.0N).
\item[(vi)] \textit{Digital twin discrepancy.} The sponsor must report any instance in which a digital twin model used for treatment planning or intraoperative guidance produces predictions that diverge from observed clinical outcomes by more than the pre-specified accuracy thresholds documented in the protocol.
\end{description}

(h) \textit{Physical AI safety review.} In addition to the review requirements of paragraph (b), the sponsor must establish a Physical AI Safety Review Committee (or equivalent function) responsible for the ongoing review and evaluation of Physical AI adverse events. The committee shall include members with expertise in robotic systems engineering, AI/ML, clinical oncology, and cybersecurity. The committee shall meet at intervals not exceeding 90 days during active clinical investigations to review Physical AI safety data and to recommend modifications to Physical AI system operation, human oversight requirements, or USL threshold criteria.

(i) \textit{Hash-chained audit trail preservation.} For all Physical AI adverse events reported under this section, the sponsor must preserve the complete hash-chained audit trail for the Physical AI system covering the period from 24 hours before the event to 72 hours after the event. The audit trail must be maintained in accordance with the electronic records requirements of 21 CFR Part 11 and must be made available to FDA upon request.

[75 FR 59961, Sept. 29, 2010. Physical AI adaptation added 17 March 2026.]

\subsection{\S~312.33 Annual Reports}

A sponsor shall within 60 days of the anniversary date that the IND went into effect, submit a brief report of the progress of the investigation that includes:

(a) \textit{Individual study information.} A brief summary of the status of each study in progress and each study completed during the previous year, including: (1) the title of the study, its purpose, a brief statement identifying the patient population, and whether the study is completed; (2) the total number of subjects initially planned, the number entered to date tabulated by age group, gender, and race, the number whose participation was completed as planned, and the number who dropped out for any reason; and (3) if the study has been completed or interim results are known, a brief description of any available study results.

(b) \textit{Summary information.} Information obtained during the previous year's clinical and nonclinical investigations, including: (1) a narrative or tabular summary of the most frequent and most serious adverse experiences by body system; (2) a summary of all IND safety reports submitted during the past year; (3) a list of subjects who died during participation in the investigation, with the cause of death for each; (4) a list of subjects who dropped out in association with any adverse experience; (5) a brief description of pertinent findings regarding the drug's actions; (6) a list of preclinical studies completed or in progress during the past year; and (7) a summary of any significant manufacturing or microbiological changes made during the past year.

(c) A description of the general investigational plan for the coming year to replace that submitted 1 year earlier.

(d) If the investigator brochure has been revised, a description of the revision and a copy of the new brochure.

(e) A description of any significant Phase 1 protocol modifications made during the previous year and not previously reported to the IND in a protocol amendment.

(f) A brief summary of significant foreign marketing developments with the drug during the past year.

(g) If desired by the sponsor, a log of any outstanding business with respect to the IND for which the sponsor requests or expects a reply, comment, or meeting.

\subsubsection{Physical AI Adaptation of \S~312.33}

(h) \textit{Physical AI system performance summary.} In addition to the information required by paragraphs (a) through (g), the annual report for investigations involving Physical AI systems shall include:

\begin{description}[leftmargin=0.8cm,style=sameline]
\item[(i)] A summary of all Physical AI adverse events reported during the previous year, organized by system type, event category (malfunction, AI error, cybersecurity, digital twin discrepancy), and patient outcome;
\item[(ii)] Updated USL scores for each Physical AI system used in the investigation, with a comparison to the scores reported in the initial IND submission or the most recent annual report, and an explanation of any changes;
\item[(iii)] A summary of AI/ML model updates implemented during the year, distinguishing between updates within PCCP scope and updates requiring protocol amendments, and including performance metrics before and after each update;
\item[(iv)] Updated simulation validation results, including any new cross-framework validation studies, updated sim-to-real gap measurements, and the results of any re-validation prompted by Physical AI system modifications;
\item[(v)] A summary of digital twin model performance, including calibration accuracy, prediction reliability, and any instances of clinically significant divergence from observed outcomes;
\item[(vi)] A summary of cybersecurity events affecting Physical AI systems, including vulnerability disclosures, patch deployments, and incident investigations;
\item[(vii)] A summary of the Physical AI Safety Review Committee's meetings, findings, and recommendations during the year; and
\item[(viii)] An updated Physical AI investigational plan for the coming year, including planned system upgrades, new Physical AI system introductions, and changes to autonomy levels or human oversight requirements.
\end{description}

[52 FR 8831, Mar. 19, 1987, as amended at 52 FR 23031, June 17, 1987; 63 FR 6862, Feb. 11, 1998; 67 FR 9585, Mar. 4, 2002. Physical AI adaptation added 17 March 2026.]

\subsection{\S~312.38 Withdrawal of an IND}

(a) At any time a sponsor may withdraw an effective IND without prejudice.

(b) If an IND is withdrawn, FDA shall be so notified, all clinical investigations conducted under the IND shall be ended, all current investigators notified, and all stocks of the drug returned to the sponsor or otherwise disposed of at the request of the sponsor in accordance with \S~312.59.

(c) If an IND is withdrawn because of a safety reason, the sponsor shall promptly so inform FDA, all participating investigators, and all reviewing Institutional Review Boards, together with the reasons for such withdrawal.

\subsubsection{Physical AI Adaptation of \S~312.38}

(d) \textit{Physical AI system decommissioning upon IND withdrawal.} When an IND is withdrawn and the investigation involved Physical AI systems, the sponsor shall, in addition to the requirements of paragraphs (a) through (c):

\begin{description}[leftmargin=0.8cm,style=sameline]
\item[(i)] Ensure that all Physical AI systems are safely deactivated and removed from clinical service in accordance with the decommissioning procedures specified in the Physical AI System Description;
\item[(ii)] Preserve all Physical AI system logs, telemetry data, hash-chained audit trails, simulation validation records, and digital twin models for the retention period specified in \S~312.57;
\item[(iii)] Securely erase all patient-specific data from Physical AI system local storage in compliance with HIPAA Safe Harbor de-identification requirements, while maintaining de-identified copies in the sponsor's central records;
\item[(iv)] Notify all sites of the Physical AI system decommissioning requirements and verify completion of decommissioning activities at each site within 30 days of the IND withdrawal; and
\item[(v)] If the withdrawal is for a safety reason related to the Physical AI system, provide FDA with a detailed technical analysis of the Physical AI system failure modes and the contributing factors, within 60 days of the withdrawal.
\end{description}

[52 FR 8831, Mar. 19, 1987, as amended at 52 FR 23031, June 17, 1987; 67 FR 9586, Mar. 4, 2002. Physical AI adaptation added 17 March 2026.]

\clearpage

\section{SUBPART C --- ADMINISTRATIVE ACTIONS}

\subsection{\S~312.40 General Requirements for Use of an Investigational New Drug in a Clinical Investigation}

(a) An investigational new drug may be used in a clinical investigation if the following conditions are met:

\begin{description}[leftmargin=0.5cm,style=sameline]
\item[(1)] The sponsor of the investigation submits an IND for the drug to FDA; the IND is in effect under paragraph (b) of this section; and the sponsor complies with all applicable requirements in this part and parts 50 and 56 with respect to the conduct of the clinical investigations; and
\item[(2)] Each participating investigator conducts his or her investigation in compliance with the requirements of this part and parts 50 and 56.
\end{description}

(b) An IND goes into effect:

\begin{description}[leftmargin=0.5cm,style=sameline]
\item[(1)] Thirty days after FDA receives the IND, unless FDA notifies the sponsor that the investigations described in the IND are subject to a clinical hold under \S~312.42; or
\item[(2)] On earlier notification by FDA that the clinical investigations in the IND may begin. FDA will notify the sponsor in writing of the date it receives the IND.
\end{description}

(c) A sponsor may ship an investigational new drug to investigators named in the IND:

\begin{description}[leftmargin=0.5cm,style=sameline]
\item[(1)] Thirty days after FDA receives the IND; or
\item[(2)] On earlier FDA authorization to ship the drug.
\end{description}

(d) An investigator may not administer an investigational new drug to human subjects until the IND goes into effect under paragraph (b) of this section.

\subsubsection{Physical AI Adaptation of \S~312.40}

(e) \textit{Physical AI system readiness requirements.} In addition to the conditions in paragraph (a), when a clinical investigation involves Physical AI systems, the following conditions must also be met before the investigational drug may be administered through or monitored by a Physical AI system:

\begin{description}[leftmargin=0.8cm,style=sameline]
\item[(i)] The Physical AI System Description required under \S~312.23(g) has been submitted as part of the IND and has not been the subject of a clinical hold under \S~312.42;
\item[(ii)] Each Physical AI system to be used in the investigation has met or exceeded the minimum USL threshold applicable to its procedure type, as specified in \S~312.1(e), or has been granted a USL threshold waiver under \S~312.10(d);
\item[(iii)] The pre-procedure safety matrix for each Physical AI system has been verified at the investigational site, including valid authorization tokens, robot readiness verification, simulation validation confirmation, and environmental checks;
\item[(iv)] All investigators and subinvestigators who will operate or supervise the Physical AI system have completed the training requirements specified in the Physical AI System Description and have demonstrated competency through assessments documented in the investigator's file; and
\item[(v)] The site's MCP server infrastructure, if applicable, has been configured and tested to the conformance level required by the protocol.
\end{description}

(f) \textit{Physical AI system deployment.} A sponsor may deploy Physical AI system components (hardware, software, simulation environments) to investigational sites named in the IND under the same timelines as paragraph (c) of this section for drug shipment. Physical AI system components shall not be used in clinical procedures involving investigational drugs until the IND is in effect and all conditions of paragraph (e) have been verified at the site.

[52 FR 8831, Mar. 19, 1987. Physical AI adaptation added 17 March 2026.]

\subsection{\S~312.41 Comment and Advice on an IND}

(a) FDA may at any time during the course of the investigation communicate with the sponsor orally or in writing about deficiencies in the IND or about FDA's need for more data or information.

(b) On the sponsor's request, FDA will provide advice on specific matters relating to an IND. Examples of such advice may include advice on the adequacy of technical data to support an investigational plan, on the design of a clinical trial, and on whether proposed investigations are likely to produce the data and information that is needed to meet requirements for a marketing application.

(c) Unless the communication is accompanied by a clinical hold order under \S~312.42, FDA communications with a sponsor under this section are solely advisory and do not require any modification in the planned or ongoing clinical investigations or response to the agency.

\subsubsection{Physical AI Adaptation of \S~312.41}

(d) \textit{Physical AI advisory communications.} FDA may communicate with the sponsor regarding deficiencies in the Physical AI System Description, simulation validation data, cybersecurity assessment, or other Physical AI-specific components of the IND. The sponsor may request FDA advice on the adequacy of Physical AI system documentation, the appropriateness of USL thresholds for specific clinical procedures, the acceptability of simulation validation methodologies, and the design of human oversight protocols for Physical AI-assisted drug administration.

(e) \textit{Physical AI system demonstrations.} At FDA's request or the sponsor's initiative, the sponsor may arrange for a demonstration of the Physical AI system's capabilities in a non-clinical setting as part of the advisory process. Such demonstrations do not constitute clinical investigations and are not subject to the IND requirements of this part. Demonstration protocols, results, and any FDA feedback shall be documented and maintained in the sponsor's IND file.

[52 FR 8831, Mar. 19, 1987, as amended at 52 FR 23031, June 17, 1987; 67 FR 9586, Mar. 4, 2002. Physical AI adaptation added 17 March 2026.]

\subsection{\S~312.42 Clinical Holds and Requests for Modification}

(a) \textit{General.} A clinical hold is an order issued by FDA to the sponsor to delay a proposed clinical investigation or to suspend an ongoing investigation. The clinical hold order may apply to one or more of the investigations covered by an IND. When a proposed study is placed on clinical hold, subjects may not be given the investigational drug. When an ongoing study is placed on clinical hold, no new subjects may be recruited to the study and placed on the investigational drug; patients already in the study should be taken off therapy involving the investigational drug unless specifically permitted by FDA in the interest of patient safety.

(b) \textit{Grounds for imposition of clinical hold.}

\begin{description}[leftmargin=0.5cm,style=sameline]
\item[(1)] \textit{Clinical hold of a Phase 1 study under an IND.} FDA may place a proposed or ongoing Phase 1 investigation on clinical hold if it finds that:
\begin{description}[leftmargin=0.8cm,style=sameline]
\item[(i)] Human subjects are or would be exposed to an unreasonable and significant risk of illness or injury;
\item[(ii)] The clinical investigators named in the IND are not qualified by reason of their scientific training and experience to conduct the investigation described in the IND;
\item[(iii)] The investigator brochure is misleading, erroneous, or materially incomplete; or
\item[(iv)] The IND does not contain sufficient information required under \S~312.23 to assess the risks to subjects of the proposed studies.
\item[(v)] The IND is for the study of an investigational drug intended to treat a life-threatening disease or condition that affects both genders, and men or women with reproductive potential who have the disease or condition being studied are excluded from eligibility because of a risk or potential risk from use of the investigational drug of reproductive toxicity or developmental toxicity, except under special circumstances as described in the regulation.
\end{description}

\item[(2)] \textit{Clinical hold of a Phase 2 or 3 study under an IND.} FDA may place a proposed or ongoing Phase 2 or 3 investigation on clinical hold if it finds that: (i) any of the conditions in paragraphs (b)(1)(i) through (b)(1)(v) of this section apply; or (ii) the plan or protocol for the investigation is clearly deficient in design to meet its stated objectives.

\item[(3)] \textit{Clinical hold of an expanded access IND or expanded access protocol.} FDA may place an expanded access IND or expanded access protocol on clinical hold if the pertinent criteria in subpart I of this part for permitting the expanded access use to begin are not satisfied, the expanded access submission does not comply with the requirements for expanded access submissions in subpart I, or the pertinent criteria for permitting the expanded access are no longer satisfied.

\item[(4)] \textit{Clinical hold of any study that is not designed to be adequate and well-controlled.} FDA may place a proposed or ongoing investigation that is not designed to be adequate and well-controlled on clinical hold if it finds that any of the conditions in paragraph (b)(1) or (b)(2) apply, or if there is reasonable evidence that the investigation is impeding enrollment in or interfering with an adequate and well-controlled study, or if insufficient quantities of the drug exist, or if the drug has been studied in adequate and well-controlled investigations that strongly suggest lack of effectiveness, among other grounds specified in the regulation.

\item[(5)] \textit{Clinical hold of any investigation involving an exception from informed consent under \S~50.24 of this chapter.} FDA may place such an investigation on clinical hold if any of the conditions in paragraphs (b)(1) or (b)(2) apply, or if the pertinent criteria in \S~50.24 are not submitted or not satisfied.

\item[(6)] \textit{Clinical hold of any investigation involving an exception from informed consent under \S~50.23(d) of this chapter.} FDA may place such an investigation on clinical hold if any of the conditions in paragraphs (b)(1) or (b)(2) apply, or if a determination by the President to waive the prior consent requirement has not been made.
\end{description}

(c) \textit{Discussion of deficiency.} Whenever FDA concludes that a deficiency exists in a clinical investigation that may be grounds for the imposition of clinical hold, FDA will, unless patients are exposed to immediate and serious risk, attempt to discuss and satisfactorily resolve the matter with the sponsor before issuing the clinical hold order.

(d) \textit{Imposition of clinical hold.} The clinical hold order may be made by telephone or other means of rapid communication or in writing. The clinical hold order will identify the studies under the IND to which the hold applies, and will briefly explain the basis for the action. As soon as possible, and no more than 30 days after imposition, the Division Director will provide the sponsor a written explanation of the basis for the hold.

(e) \textit{Resumption of clinical investigations.} An investigation may only resume after FDA has notified the sponsor that the investigation may proceed. Resumption will be authorized when the sponsor corrects the deficiency(ies) previously cited or otherwise satisfies the agency that the investigation(s) can proceed. If a sponsor requests in writing that the clinical hold be removed and submits a complete response, FDA shall respond in writing within 30 calendar days.

(f) \textit{Appeal.} If the sponsor disagrees with the reasons cited for the clinical hold, the sponsor may request reconsideration in accordance with \S~312.48.

(g) \textit{Conversion of IND on clinical hold to inactive status.} If all investigations covered by an IND remain on clinical hold for 1 year or more, the IND may be placed on inactive status by FDA under \S~312.45.

\subsubsection{Physical AI Adaptation of \S~312.42}

(h) \textit{Physical AI grounds for clinical hold.} In addition to the grounds specified in paragraphs (b)(1) through (b)(6), FDA may place a proposed or ongoing clinical investigation on clinical hold based on Physical AI system deficiencies, including:

\begin{description}[leftmargin=0.8cm,style=sameline]
\item[(i)] \textit{Robotic safety failure.} The Physical AI system has demonstrated a pattern of safety-critical failures, including unplanned emergency stop activations exceeding the rate specified in the protocol, collision events resulting in or potentially resulting in patient injury, or force limit exceedances during clinical procedures;
\item[(ii)] \textit{AI model degradation.} One or more AI/ML models embedded in the Physical AI system have demonstrated sustained performance degradation below the pre-specified acceptance criteria, and the sponsor has not implemented effective corrective action within the timeframe specified in the PCCP or Physical AI System Description;
\item[(iii)] \textit{Simulation-reality divergence.} The observed behavior of the Physical AI system during clinical procedures diverges from the validated simulation predictions by more than twice the established sim-to-real gap tolerance, indicating that the pre-clinical simulation validation no longer adequately predicts the system's clinical behavior;
\item[(iv)] \textit{Cybersecurity compromise.} The Physical AI system has been the subject of a confirmed cybersecurity breach that compromises or has the potential to compromise patient safety, data integrity, or the security of the investigational drug administration pathway;
\item[(v)] \textit{USL score decline.} The Physical AI system's USL score has declined below the minimum threshold required for its procedure type, whether due to identified deficiencies in simulation switching, AI integration, cross-robot sharing, or clinical trial collaboration capabilities;
\item[(vi)] \textit{Inadequate Physical AI System Description.} The Physical AI System Description submitted under \S~312.23(g) does not contain sufficient information to assess the risks to subjects posed by the Physical AI system, including inadequate simulation validation data, incomplete cybersecurity assessment, or insufficient human oversight protocols;
\item[(vii)] \textit{Digital twin failure.} A digital twin model used for treatment planning or intraoperative guidance has produced clinically significant erroneous predictions that resulted in or could have resulted in inappropriate drug administration, and the sponsor has not demonstrated adequate corrective action; and
\item[(viii)] \textit{Human oversight failure.} The investigation demonstrates a pattern of inadequate human oversight of Physical AI systems, including instances where operators failed to intervene when system performance indicators warranted intervention, or where the operator-to-system ratio was insufficient to maintain effective supervisory control.
\end{description}

(i) \textit{Physical AI clinical hold procedures.} When a clinical hold is imposed on the basis of Physical AI system deficiencies under paragraph (h):

\begin{description}[leftmargin=0.8cm,style=sameline]
\item[(i)] The sponsor shall immediately cease all Physical AI system operations in clinical procedures involving the investigational drug at all affected sites. Patients currently undergoing a Physical AI-assisted procedure at the time of the hold order shall have the procedure completed using manual methods under the supervising investigator's clinical judgment, unless FDA specifically authorizes completion of the procedure using the Physical AI system in the interest of patient safety;
\item[(ii)] The sponsor shall place all affected Physical AI systems in a safe standby mode and preserve all system logs, telemetry data, and hash-chained audit trails from the period preceding the clinical hold;
\item[(iii)] The clinical hold order shall specify whether the hold applies to all Physical AI systems under the IND or only to specific system types, models, or sites; and
\item[(iv)] To resume Physical AI system operations, the sponsor must demonstrate that the Physical AI deficiency has been corrected and that the system has been re-validated through the simulation validation process described in \S~312.21(d), unless FDA determines that a lesser level of re-validation is sufficient based on the nature of the deficiency.
\end{description}

[52 FR 8831, Mar. 19, 1987, as amended at 52 FR 19477, May 22, 1987; 57 FR 13249, Apr. 15, 1992; 61 FR 51530, Oct. 2, 1996; 63 FR 68678, Dec. 14, 1998; 64 FR 54189, Oct. 5, 1999; 65 FR 34971, June 1, 2000; 74 FR 40942, Aug. 13, 2009. Physical AI adaptation added 17 March 2026.]

\subsection{\S~312.44 Termination}

(a) \textit{General.} This section describes the procedures under which FDA may terminate an IND. If an IND is terminated, the sponsor shall end all clinical investigations conducted under the IND and recall or otherwise provide for the disposition of all unused supplies of the drug. A termination action may be based on deficiencies in the IND or in the conduct of an investigation under an IND. Except as provided in paragraph (d) of this section, a termination shall be preceded by a proposal to terminate by FDA and an opportunity for the sponsor to respond. FDA will, in general, only initiate an action under this section after first attempting to resolve differences informally or through the clinical hold procedures described in \S~312.42.

(b) \textit{Grounds for termination.}

\begin{description}[leftmargin=0.5cm,style=sameline]
\item[(1)] \textit{Phase 1.} FDA may propose to terminate an IND during Phase 1 if it finds that:
\begin{description}[leftmargin=0.8cm,style=sameline]
\item[(i)] Human subjects would be exposed to an unreasonable and significant risk of illness or injury.
\item[(ii)] The IND does not contain sufficient information required under \S~312.23 to assess the safety to subjects.
\item[(iii)] The methods, facilities, and controls used for manufacturing, processing, and packing of the investigational drug are inadequate.
\item[(iv)] The clinical investigations are being conducted in a manner substantially different than that described in the protocols submitted in the IND.
\item[(v)] The drug is being promoted or distributed for commercial purposes not justified by the requirements of the investigation.
\item[(vi)] The IND contains an untrue statement of a material fact or omits material information required by this part.
\item[(vii)] The sponsor fails promptly to investigate and inform FDA and all investigators of serious and unexpected adverse experiences in accordance with \S~312.32 or fails to make any other report required under this part.
\item[(viii)] The sponsor fails to submit an accurate annual report in accordance with \S~312.33.
\item[(ix)] The sponsor fails to comply with any other applicable requirement of this part, part 50, or part 56.
\item[(x)] The IND has remained on inactive status for 5 years or more.
\item[(xi)] The sponsor fails to delay a proposed investigation or to suspend an ongoing investigation that has been placed on clinical hold under \S~312.42(b)(4).
\end{description}

\item[(2)] \textit{Phase 2 or 3.} FDA may propose to terminate an IND during Phase 2 or Phase 3 if FDA finds that: (i) any of the conditions in paragraphs (b)(1)(i) through (b)(1)(xi) apply; or (ii) the investigational plan or protocol(s) is not reasonable as a bona fide scientific plan to determine whether or not the drug is safe and effective for use; or (iii) there is convincing evidence that the drug is not effective for the purpose for which it is being investigated.

\item[(3)] FDA may propose to terminate a treatment IND if it finds that: (i) any of the conditions in paragraphs (b)(1)(i) through (x) apply; or (ii) any of the conditions in \S~312.42(b)(3) apply.
\end{description}

(c) \textit{Opportunity for sponsor response.} If FDA proposes to terminate an IND, FDA will notify the sponsor in writing and invite correction or explanation within 30 days. If the sponsor does not respond, the IND shall be terminated. If the sponsor responds but FDA does not accept the explanation, the sponsor shall be informed and provided an opportunity for a regulatory hearing under part 16.

(d) \textit{Immediate termination of IND.} Notwithstanding paragraphs (a) through (c), if at any time FDA concludes that continuation of the investigation presents an immediate and substantial danger to the health of individuals, the agency shall immediately terminate the IND by written notice to the sponsor from the Director of the Center for Drug Evaluation and Research or the Director of the Center for Biologics Evaluation and Research. An IND so terminated is subject to reinstatement by the Director on the basis of additional submissions that eliminate such danger.

\subsubsection{Physical AI Adaptation of \S~312.44}

(e) \textit{Physical AI grounds for termination.} In addition to the grounds specified in paragraph (b), FDA may propose to terminate an IND based on Physical AI system deficiencies that present an ongoing and unresolvable risk to human subjects, including:

\begin{description}[leftmargin=0.8cm,style=sameline]
\item[(i)] Repeated Physical AI safety failures that have not been corrected despite clinical hold and sponsor response, demonstrating a fundamental inadequacy in the Physical AI system's design, manufacture, or operation for the intended clinical use;
\item[(ii)] A confirmed cybersecurity compromise of the Physical AI system that results in unauthorized modification of drug administration parameters, patient data exposure, or loss of system integrity, and for which the sponsor cannot demonstrate adequate remediation;
\item[(iii)] Systematic failure of the sponsor to report Physical AI adverse events in accordance with \S~312.32(g), to maintain Physical AI system logs and audit trails in accordance with \S~312.57, or to comply with the Physical AI-specific requirements of this part; and
\item[(iv)] Evidence that the Physical AI system's simulation validation was materially deficient, misleading, or falsified, such that the pre-clinical safety basis for the investigation is fundamentally compromised.
\end{description}

(f) \textit{Immediate termination for Physical AI emergencies.} FDA may immediately terminate an IND under paragraph (d) when a Physical AI system malfunction presents an immediate and substantial danger to patient health. Such circumstances include a Physical AI system that has caused or is likely to cause serious patient injury due to uncontrolled robotic movement, incorrect drug delivery at dangerous doses, or loss of emergency stop functionality. Upon immediate termination, the sponsor shall decommission all Physical AI systems in accordance with \S~312.38(d).

[52 FR 8831, Mar. 19, 1987, as amended at 52 FR 23031, June 17, 1987; 55 FR 11579, Mar. 29, 1990; 57 FR 13249, Apr. 15, 1992; 67 FR 9586, Mar. 4, 2002. Physical AI adaptation added 17 March 2026.]

\subsection{\S~312.45 Inactive Status}

(a) If no subjects are entered into clinical studies for a period of 2 years or more under an IND, or if all investigations under an IND remain on clinical hold for 1 year or more, the IND may be placed by FDA on inactive status. This action may be taken by FDA either on request of the sponsor or on FDA's own initiative. If FDA seeks to act on its own initiative, it shall first notify the sponsor in writing, and the sponsor shall have 30 days to respond as to why the IND should continue to remain active.

(b) If an IND is placed on inactive status, all investigators shall be so notified and all stocks of the drug shall be returned or otherwise disposed of in accordance with \S~312.59.

(c) A sponsor is not required to submit annual reports to an IND on inactive status. An inactive IND is, however, still in effect for purposes of the public disclosure of data and information under \S~312.130.

(d) A sponsor who intends to resume clinical investigation under an IND placed on inactive status shall submit a protocol amendment under \S~312.30 containing the proposed general investigational plan and appropriate protocols. Clinical investigations under an IND on inactive status may only resume 30 days after FDA receives the protocol amendment, unless FDA notifies the sponsor that the investigations are subject to a clinical hold, or on earlier notification by FDA.

(e) An IND that remains on inactive status for 5 years or more may be terminated under \S~312.44.

\subsubsection{Physical AI Adaptation of \S~312.45}

(f) \textit{Physical AI system dormancy.} When an IND is placed on inactive status and the investigation involved Physical AI systems, the sponsor shall:

\begin{description}[leftmargin=0.8cm,style=sameline]
\item[(i)] Place all Physical AI systems in a documented dormancy state, which includes powering down robotic hardware, archiving simulation validation data, and securing all system credentials and access tokens;
\item[(ii)] Maintain Physical AI system maintenance records, calibration data, and software version documentation sufficient to support system reactivation; and
\item[(iii)] Continue cybersecurity monitoring of any Physical AI system components that remain connected to clinical networks during the dormancy period, including patch management for critical security vulnerabilities.
\end{description}

(g) \textit{Physical AI system reactivation.} When a sponsor seeks to resume clinical investigation under an IND that was placed on inactive status and involved Physical AI systems, the protocol amendment required under paragraph (d) shall include, in addition to the standard requirements:

\begin{description}[leftmargin=0.8cm,style=sameline]
\item[(i)] Updated Physical AI System Description reflecting the current hardware configuration, software versions, and AI/ML model versions;
\item[(ii)] Current USL scores for each Physical AI system, reassessed using the most recent version of the USL scoring methodology;
\item[(iii)] Updated simulation validation data demonstrating that the Physical AI system continues to meet cross-framework behavioral consistency requirements after the dormancy period;
\item[(iv)] Updated cybersecurity assessment addressing any vulnerabilities identified during the dormancy period and confirming that all security patches have been applied; and
\item[(v)] Evidence that all investigators and subinvestigators have completed refresher training on the Physical AI system and have demonstrated current competency.
\end{description}

[52 FR 8831, Mar. 19, 1987, as amended at 52 FR 23031, June 17, 1987; 67 FR 9586, Mar. 4, 2002. Physical AI adaptation added 17 March 2026.]

\subsection{\S~312.47 Meetings}

(a) \textit{General.} Meetings between a sponsor and the agency are frequently useful in resolving questions and issues raised during the course of a clinical investigation. FDA encourages such meetings to the extent that they aid in the evaluation of the drug and in the solution of scientific problems concerning the drug, to the extent that FDA's resources permit. The general principle underlying the conduct of such meetings is that there should be free, full, and open communication about any scientific or medical question that may arise during the clinical investigation. These meetings shall be conducted and documented in accordance with part 10.

(b) \textit{``End-of-Phase 2'' meetings and meetings held before submission of a marketing application.}

\begin{description}[leftmargin=0.5cm,style=sameline]
\item[(1)] \textit{End-of-Phase 2 meetings.} The purpose of an end-of-phase 2 meeting is to determine the safety of proceeding to Phase 3, to evaluate the Phase 3 plan and protocols and the adequacy of current studies and plans to assess pediatric safety and effectiveness, and to identify any additional information necessary to support a marketing application. While designed primarily for INDs involving new molecular entities or major new uses, a sponsor of any IND may request and obtain an end-of-Phase 2 meeting. The sponsor should submit background information at least 1 month in advance, including summaries of Phase 1 and 2 investigations, specific Phase 3 protocols, plans for additional nonclinical studies, pediatric study plans, and tentative labeling if available. Agreements reached at the meeting will be recorded in minutes taken by FDA and provided to the sponsor.

\item[(2)] \textit{``Pre-NDA'' and ``pre-BLA'' meetings.} The primary purpose of these meetings is to uncover any major unresolved problems, to identify adequate and well-controlled studies establishing effectiveness, to identify the status of ongoing pediatric studies, to acquaint FDA reviewers with the submission, and to discuss statistical analysis methods and data presentation. The sponsor should submit background information at least 1 month in advance, including a brief summary of clinical studies, a proposed format, information on pediatric studies, and any other information for discussion.
\end{description}

\subsubsection{Physical AI Adaptation of \S~312.47}

(c) \textit{Physical AI system review meetings.} In addition to the meetings described in paragraphs (a) and (b), the following Physical AI-specific meetings are available:

\begin{description}[leftmargin=0.8cm,style=sameline]
\item[(i)] \textit{Pre-IND Physical AI system review.} Prior to submitting an IND that includes a Physical AI system, the sponsor may request a meeting with FDA to review the Physical AI System Description, simulation validation approach, cybersecurity assessment methodology, and proposed human oversight protocols. The sponsor should submit the draft Physical AI System Description at least 1 month in advance.
\item[(ii)] \textit{End-of-Phase 2 Physical AI review.} When an end-of-Phase 2 meeting is requested for an investigation involving Physical AI systems, the sponsor should include in the advance submission materials a summary of Physical AI system performance during Phases 1 and 2, including Physical AI adverse event rates, system reliability metrics, USL score trends, and the proposed Physical AI system configuration for Phase 3 multi-site deployment.
\item[(iii)] \textit{Physical AI system demonstration.} At any meeting described in this section, the sponsor may offer to provide a non-clinical demonstration of the Physical AI system's capabilities, including simulated procedure execution, emergency stop activation, and human-robot interaction workflows. Such demonstrations shall be conducted in a controlled non-clinical environment and shall be documented in the meeting minutes.
\end{description}

[52 FR 8831, Mar. 19, 1987, as amended at 52 FR 23031, June 17, 1987; 55 FR 11580, Mar. 29, 1990; 63 FR 66669, Dec. 2, 1998; 67 FR 9586, Mar. 4, 2002. Physical AI adaptation added 17 March 2026.]

\subsection{\S~312.48 Dispute Resolution}

(a) \textit{General.} The Food and Drug Administration is committed to resolving differences between sponsors and FDA reviewing divisions with respect to requirements for IND's as quickly and amicably as possible through the cooperative exchange of information and views.

(b) \textit{Administrative and procedural issues.} When administrative or procedural disputes arise, the sponsor should first attempt to resolve the matter with the division in FDA's Center for Drug Evaluation and Research or Center for Biologics Evaluation and Research which is responsible for review of the IND, beginning with the consumer safety officer assigned to the application. If the dispute is not resolved, the sponsor may raise the matter with the person designated as ombudsman, whose function shall be to investigate what has happened and to facilitate a timely and equitable resolution.

(c) \textit{Scientific and medical disputes.}

\begin{description}[leftmargin=0.5cm,style=sameline]
\item[(1)] When scientific or medical disputes arise during the drug investigation process, sponsors should discuss the matter directly with the responsible reviewing officials. If necessary, sponsors may request a meeting with the appropriate reviewing officials and management representatives in order to seek a resolution.
\item[(2)] The ``end-of-Phase 2'' and ``pre-NDA'' meetings described in \S~312.47(b) will also provide a timely forum for discussing and resolving scientific and medical issues.
\item[(3)] In requesting a meeting designed to resolve a scientific or medical dispute, applicants may suggest that FDA seek the advice of outside experts, in which case FDA may invite to the meeting one or more advisory committee members or other consultants. For major scientific and medical policy issues not resolved by informal meetings, FDA may refer the matter to a standing advisory committee.
\end{description}

\subsubsection{Physical AI Adaptation of \S~312.48}

(d) \textit{Physical AI technical disputes.} When disputes arise regarding Physical AI system requirements, including USL threshold interpretations, simulation validation adequacy, cybersecurity assessment sufficiency, or the appropriateness of human oversight protocols, the following provisions apply in addition to paragraphs (a) through (c):

\begin{description}[leftmargin=0.8cm,style=sameline]
\item[(i)] \textit{Technical review panel.} For disputes involving Physical AI system safety or performance, the sponsor may request that FDA convene a technical review panel that includes members with expertise in robotic systems engineering, AI/ML validation, simulation physics, cybersecurity, and clinical oncology, in addition to the standard reviewing officials;
\item[(ii)] \textit{Independent simulation validation.} When a dispute centers on the adequacy of simulation validation data, either party may propose that an independent simulation validation study be conducted using the physical-ai-oncology-trials validation benchmark suite (v1.0) or an equivalent validated benchmark, with the results to be reviewed by both parties; and
\item[(iii)] \textit{USL score disputes.} When a dispute centers on a Physical AI system's USL score or the applicable USL threshold for a specific clinical procedure, the sponsor may submit an independent USL assessment conducted by a qualified third party, which FDA shall consider in its review.
\end{description}

[52 FR 8831, Mar. 19, 1987, as amended at 55 FR 11580, Mar. 29, 1990. Physical AI adaptation added 17 March 2026.]

\clearpage

\section{SUBPART D --- RESPONSIBILITIES OF SPONSORS AND INVESTIGATORS}

\subsection{\S~312.50 General Responsibilities of Sponsors}

A sponsor is responsible for selecting qualified investigators, providing them with the information they need to conduct an investigation properly, ensuring proper monitoring of the investigation(s), ensuring that the investigation(s) is conducted in accordance with the general investigational plan and protocols contained in the IND, maintaining an effective IND with respect to the investigations, and ensuring that FDA and all participating investigators are promptly informed of significant new adverse effects or risks with respect to the drug. Additional specific responsibilities of sponsors are described elsewhere in this part.

\subsubsection{Physical AI Adaptation of \S~312.50}

In addition to the general responsibilities described above, a sponsor of a clinical investigation that involves Physical AI systems shall be responsible for:

\begin{description}[leftmargin=0.8cm,style=sameline]
\item[(a)] \textit{Physical AI system selection and qualification.} Selecting Physical AI systems that are appropriate for the intended clinical procedures and that meet or exceed the minimum USL thresholds specified in \S~312.1(e). The sponsor shall maintain documentation of the selection rationale, including comparative analysis of available Physical AI platforms and justification for the selected system's suitability for oncology drug administration;
\item[(b)] \textit{Investigator Physical AI training.} Ensuring that all investigators and subinvestigators who will operate or supervise Physical AI systems have completed the training specified in the Physical AI System Description and have demonstrated competency through documented assessments;
\item[(c)] \textit{Physical AI system deployment and maintenance.} Ensuring that Physical AI systems are properly deployed, configured, calibrated, and maintained at each investigational site in accordance with the manufacturer's specifications and the Physical AI System Description;
\item[(d)] \textit{Physical AI monitoring.} Ensuring that the Physical AI system performance is monitored continuously during clinical procedures and that monitoring data, including telemetry, safety events, and performance metrics, are captured and retained in accordance with the data management plan;
\item[(e)] \textit{Physical AI adverse event reporting.} Ensuring that Physical AI adverse events are identified, documented, and reported to FDA and all participating investigators in accordance with the Physical AI adverse event reporting requirements of \S~312.32(g);
\item[(f)] \textit{Cybersecurity management.} Maintaining the cybersecurity posture of all Physical AI systems throughout the investigation, including timely application of security patches, monitoring for unauthorized access, and incident response in accordance with the cybersecurity assessment submitted under \S~312.23(g)(iv); and
\item[(g)] \textit{Physical AI system updates.} Managing all software updates, AI/ML model updates, and hardware modifications to the Physical AI system in accordance with the PCCP and the amendment requirements of \S~312.30.
\end{description}

[52 FR 8831, Mar. 19, 1987. Physical AI adaptation added 17 March 2026.]

\subsection{\S~312.52 Transfer of Obligations to a Contract Research Organization}

(a) A sponsor may transfer responsibility for any or all of the obligations set forth in this part to a contract research organization. Any such transfer shall be described in writing. If not all obligations are transferred, the writing is required to describe each obligation being assumed by the contract research organization.

(b) A sponsor may not transfer to a contract research organization the obligation to submit an IND.

(c) A contract research organization that assumes any obligation of a sponsor shall comply with the specific regulatory requirements in this chapter applicable to this obligation and shall be subject to the same regulatory action as a sponsor for failure to comply with any obligation assumed under these regulations.

\subsubsection{Physical AI Adaptation of \S~312.52}

(d) \textit{Transfer of Physical AI obligations.} When a sponsor transfers responsibilities involving Physical AI systems to a contract research organization (CRO), the following additional requirements apply:

\begin{description}[leftmargin=0.8cm,style=sameline]
\item[(i)] The written transfer agreement shall explicitly identify which Physical AI-specific obligations are being transferred, including Physical AI system deployment and maintenance, Physical AI adverse event monitoring and reporting, cybersecurity management, simulation validation oversight, and Physical AI system performance monitoring;
\item[(ii)] The CRO shall demonstrate that it possesses or has contracted for the technical expertise necessary to fulfill the transferred Physical AI obligations, including personnel qualified in robotic systems engineering, AI/ML operations, and cybersecurity;
\item[(iii)] The sponsor retains responsibility for submitting all Physical AI-related IND amendments and for all Physical AI adverse event reports to FDA, even when day-to-day Physical AI system management has been transferred to the CRO, unless the transfer agreement explicitly provides otherwise and FDA has been notified; and
\item[(iv)] The transfer agreement shall specify procedures for the CRO to promptly notify the sponsor of any Physical AI safety events, cybersecurity incidents, or USL score changes that may require IND amendment or reporting to FDA.
\end{description}

[52 FR 8831, Mar. 19, 1987, as amended at 67 FR 9586, Mar. 4, 2002. Physical AI adaptation added 17 March 2026.]

\subsection{\S~312.53 Selecting Investigators and Monitors}

(a) \textit{Selecting investigators.} A sponsor shall select only investigators qualified by training and experience as appropriate experts to investigate the drug.

(b) \textit{Control of the drug.} A sponsor shall ship investigational new drugs only to investigators participating in the investigation.

(c) \textit{Providing information to the investigator.} Before the investigation begins, a sponsor (other than a sponsor-investigator) shall provide to each investigator a copy of the investigator's brochure containing the information described in \S~312.23(a)(5) and shall provide the investigator with the information needed to conduct the investigation properly.

(d) \textit{Selecting monitors.} A sponsor shall select a monitor qualified by training and experience to monitor the progress of the investigation.

\subsubsection{Physical AI Adaptation of \S~312.53}

(e) \textit{Physical AI investigator qualifications.} In addition to the requirements of paragraph (a), when selecting investigators for a clinical investigation involving Physical AI systems, the sponsor shall ensure that each investigator:

\begin{description}[leftmargin=0.8cm,style=sameline]
\item[(i)] Has completed the Physical AI system training program specified in the Physical AI System Description and has demonstrated competency through documented assessments, including hands-on operation of the Physical AI system in a non-clinical setting;
\item[(ii)] Has experience or training in supervising robotic-assisted medical procedures, or has completed a supervised training period using the specific Physical AI system under the guidance of an experienced operator; and
\item[(iii)] Has designated at least one qualified backup operator at the investigational site who has also completed the training requirements, to ensure continuous availability of qualified supervision during Physical AI-assisted procedures.
\end{description}

(f) \textit{Physical AI monitor qualifications.} In addition to the requirements of paragraph (d), monitors for investigations involving Physical AI systems shall have training or experience sufficient to evaluate Physical AI system performance data, interpret Physical AI adverse event reports, and assess compliance with Physical AI-specific protocol requirements.

(g) \textit{Physical AI information for investigators.} In addition to the investigator's brochure required under paragraph (c), the sponsor shall provide each investigator with:

\begin{description}[leftmargin=0.8cm,style=sameline]
\item[(i)] The Physical AI System Description or a summary thereof, including system capabilities, limitations, contraindications, and known failure modes;
\item[(ii)] Physical AI system operating manuals and quick reference guides for all procedure types authorized under the protocol;
\item[(iii)] Emergency procedures for Physical AI system malfunctions, including manual override protocols and emergency stop activation procedures; and
\item[(iv)] The Physical AI adverse event reporting requirements, including the classification of Physical AI adverse events and the timelines for reporting.
\end{description}

[52 FR 8831, Mar. 19, 1987, as amended at 52 FR 23031, June 17, 1987; 61 FR 51530, Oct. 2, 1996; 67 FR 9586, Mar. 4, 2002. Physical AI adaptation added 17 March 2026.]

\subsection{\S~312.54 Emergency Research Under \S~50.24 of This Chapter}

(a) The sponsor shall monitor the progress of all investigations involving an exception from informed consent under \S~50.24 of this chapter. When the sponsor receives from the IRB information concerning the public disclosures required by \S~50.24(a)(7)(ii) and (a)(7)(iii) of this chapter, the sponsor promptly shall submit to the IND file and to Docket Number 95S--0158 in the Division of Dockets Management (HFA--305), Food and Drug Administration, 5630 Fishers Lane, rm.~1061, Rockville, MD 20852, copies of the information that was disclosed, identified by the IND number.

(b) The sponsor also shall monitor such investigations to identify when an IRB determines that it cannot approve the research because it does not meet the criteria in the exception in \S~50.24(a) of this chapter or because of other relevant ethical concerns. The sponsor promptly shall provide this information in writing to FDA, all participating investigators, and other sponsors of investigations involving waiver of informed consent for the same drug.

\subsubsection{Physical AI Adaptation of \S~312.54}

(c) \textit{Physical AI in emergency research.} When a clinical investigation involving a Physical AI system is conducted under an exception from informed consent under \S~50.24 of this chapter, the following additional requirements apply:

\begin{description}[leftmargin=0.8cm,style=sameline]
\item[(i)] The public disclosures required under \S~50.24(a)(7)(ii) and (a)(7)(iii) shall include a description of the Physical AI system's role in the investigation, including the types of procedures the Physical AI system will perform and the human oversight protocols in place;
\item[(ii)] The community consultation process shall specifically address community concerns regarding the use of Physical AI systems in emergency medical procedures without prior informed consent; and
\item[(iii)] The protocol shall include heightened Physical AI safety monitoring requirements, including real-time human oversight with immediate manual override capability for all Physical AI-assisted procedures conducted under the emergency research exception.
\end{description}

[62 FR 32479, June 16, 1997, as amended at 67 FR 9586, Mar. 4, 2002. Physical AI adaptation added 17 March 2026.]

\subsection{\S~312.55 Informing Investigators}

(a) Before the investigation begins, a sponsor (other than a sponsor-investigator) shall provide to each investigator a copy of the investigator's brochure containing the information described in \S~312.23(a)(5).

(b) The sponsor shall, as the overall investigation proceeds, keep each participating investigator informed of new observations discovered by or reported to the sponsor on the drug, particularly with respect to adverse effects and safe use. Such information may be distributed to investigators by means of periodically revised investigator brochures, reprints or published studies, reports or letters to clinical investigators, or other appropriate means.

\subsubsection{Physical AI Adaptation of \S~312.55}

(c) \textit{Physical AI system updates to investigators.} In addition to the requirements of paragraphs (a) and (b), the sponsor shall keep each participating investigator informed of:

\begin{description}[leftmargin=0.8cm,style=sameline]
\item[(i)] All Physical AI adverse events reported at other investigational sites, including the nature of the event, the Physical AI system model and software version involved, and any corrective actions taken;
\item[(ii)] Physical AI system software updates, AI/ML model updates, and hardware modifications that may affect the system's performance or safety profile, including the rationale for the update and any re-training or re-validation requirements;
\item[(iii)] Changes in USL scores for the Physical AI system model being used in the investigation, including the reasons for any score changes and the potential impact on clinical procedures;
\item[(iv)] Cybersecurity alerts, vulnerabilities, or incidents affecting the Physical AI system platform, including recommended mitigations; and
\item[(v)] Updated simulation validation data or changes to the sim-to-real gap tolerances that may affect clinical procedure planning.
\end{description}

[52 FR 8831, Mar. 19, 1987. Physical AI adaptation added 17 March 2026.]

\subsection{\S~312.56 Review of Ongoing Investigations}

(a) A sponsor who discovers that an investigator is not complying with the signed agreement (Form FDA 1572), the general investigational plan, or the requirements of this part or other applicable parts shall promptly either secure compliance, or end the investigator's participation in the investigation.

(b) A sponsor shall review and evaluate the evidence relating to the safety and effectiveness of the drug as it is obtained from the investigator. The sponsor shall make such reports to FDA regarding information relevant to the safety of the drug as are required under \S~312.32. The sponsor shall make annual reports on the progress of the investigation in accordance with \S~312.33.

(c) A sponsor shall promptly investigate and report to FDA and to all participating investigators any adverse experience associated with use of the drug that is both serious and unexpected.

(d) A sponsor of a clinical study who has not submitted a marketing application for the drug under study shall maintain records of the composition, manufacture, and controls of the drug, and make such records available for FDA inspection upon request.

\subsubsection{Physical AI Adaptation of \S~312.56}

(e) \textit{Physical AI compliance review.} In addition to the requirements of paragraph (a), the sponsor shall review investigator compliance with Physical AI-specific requirements, including:

\begin{description}[leftmargin=0.8cm,style=sameline]
\item[(i)] Proper use of the Physical AI system in accordance with the approved protocol and Physical AI System Description;
\item[(ii)] Maintenance of Physical AI system logs, telemetry data, and audit trails as required by \S~312.62;
\item[(iii)] Timely reporting of Physical AI adverse events by investigators to the sponsor;
\item[(iv)] Adherence to the human oversight requirements specified in the protocol, including maintenance of the required operator-to-system ratio; and
\item[(v)] Proper implementation of Physical AI system updates and security patches as directed by the sponsor.
\end{description}

(f) \textit{Physical AI performance review.} The sponsor shall continuously review and evaluate Physical AI system performance data across all investigational sites, including:

\begin{description}[leftmargin=0.8cm,style=sameline]
\item[(i)] Aggregate Physical AI adverse event rates and trends;
\item[(ii)] Cross-site comparison of Physical AI system performance metrics;
\item[(iii)] USL score stability and trends;
\item[(iv)] Simulation-to-reality gap measurements; and
\item[(v)] Cybersecurity incident frequency and severity.
\end{description}

The sponsor shall include a Physical AI system performance summary in each annual report submitted under \S~312.33.

[52 FR 8831, Mar. 19, 1987, as amended at 52 FR 23031, June 17, 1987; 67 FR 9586, Mar. 4, 2002. Physical AI adaptation added 17 March 2026.]

\subsection{\S~312.57 Recordkeeping and Record Retention}

(a) A sponsor shall maintain adequate records showing the receipt, shipment, or other disposition of the investigational drug. These records are required to include, as appropriate, the name of the investigator to whom the drug is shipped, and the date, quantity, and batch or code mark of each such shipment.

(b) A sponsor shall maintain complete and accurate records showing any financial interest in accordance with part 54 of this chapter, as required by \S~312.23(a)(11), and \S~312.33(a)(2).

(c) A sponsor shall retain the records and reports required by this part for 2 years after a marketing application is approved for the drug; or, if an application is not approved for the drug, until 2 years after shipment and delivery of the drug for investigational use is discontinued and FDA has been so notified.

(d) A sponsor shall retain reserve samples of any test article and reference standard for a period of not less than 5 years after the completion or discontinuation of the investigation.

\subsubsection{Physical AI Adaptation of \S~312.57}

(e) \textit{Physical AI system records.} In addition to the records required under paragraphs (a) through (d), the sponsor shall maintain the following Physical AI-specific records:

\begin{description}[leftmargin=0.8cm,style=sameline]
\item[(i)] \textit{Physical AI system deployment records.} Complete records of all Physical AI system deployments, including the system model, serial number, hardware configuration, software version, AI/ML model version, and deployment date for each system deployed to each investigational site;
\item[(ii)] \textit{Physical AI system maintenance records.} Records of all Physical AI system maintenance activities, including calibration, repairs, hardware replacements, software updates, AI/ML model updates, and security patches, with dates, descriptions, and the identity of the personnel performing each activity;
\item[(iii)] \textit{Simulation validation records.} Complete records of all simulation validation activities, including the simulation framework used, validation parameters, test results, sim-to-real gap measurements, and any cross-framework behavioral consistency assessments;
\item[(iv)] \textit{Physical AI telemetry archives.} Archived telemetry data from all Physical AI-assisted clinical procedures, including robotic movement trajectories, force/torque measurements, sensor readings, AI decision logs, and operator interaction records. Telemetry data shall be stored in a format that supports retrospective analysis and shall be hash-chained to prevent tampering;
\item[(v)] \textit{USL assessment records.} Records of all USL assessments and reassessments for each Physical AI system, including the raw scoring data, the assessor's identity and qualifications, and the final USL score;
\item[(vi)] \textit{Cybersecurity records.} Records of all cybersecurity assessments, vulnerability scans, penetration tests, security incidents, and remediation actions related to the Physical AI systems; and
\item[(vii)] \textit{Training records.} Records of all Physical AI system training provided to investigators, subinvestigators, and other site personnel, including the training content, dates, and competency assessment results.
\end{description}

(f) \textit{Physical AI record retention.} Physical AI-specific records described in paragraph (e) shall be retained for the same period specified in paragraph (c). Physical AI telemetry archives under paragraph (e)(iv) may be stored in compressed format after the completion of the investigation, provided that the data can be decompressed and rendered in a human-readable format within a reasonable time upon FDA request.

[52 FR 8831, Mar. 19, 1987, as amended at 63 FR 5252, Feb. 2, 1998; 67 FR 9586, Mar. 4, 2002. Physical AI adaptation added 17 March 2026.]

\subsection{\S~312.58 Inspection of Sponsor's Records and Reports}

(a) FDA inspection. A sponsor shall upon request from any properly authorized officer or employee of the Food and Drug Administration, at reasonable times, permit such officer or employee to have access to and copy and verify any records and reports relating to a clinical investigation conducted under this part. Upon written request by FDA, the sponsor shall submit the records or reports (or copies of them) to FDA. The sponsor shall discontinue shipments of the drug to any investigator who has failed to maintain or make available records or reports of the investigation as required by this part.

(b) Controlled substances. If an investigational new drug is a substance listed in any schedule of the Controlled Substances Act (21 U.S.C. 801; 21 CFR part 1308), records concerning shipment, delivery, receipt, and disposition of the drug, which are required to be maintained under this part or other applicable parts of this chapter shall, upon the request of a properly authorized employee of the Drug Enforcement Administration of the Department of Justice, be made available by the investigator or sponsor to whom the request is made, for inspection and copying.

\subsubsection{Physical AI Adaptation of \S~312.58}

(c) \textit{Physical AI system inspection.} In addition to the records and reports described in paragraph (a), the sponsor shall make available to FDA for inspection:

\begin{description}[leftmargin=0.8cm,style=sameline]
\item[(i)] All Physical AI system records described in \S~312.57(e), including telemetry archives, simulation validation records, and cybersecurity records;
\item[(ii)] Access to the Physical AI system software, including source code for custom-developed components directly related to drug administration safety, AI/ML model architectures and training data documentation, and configuration files. The sponsor may request confidential treatment of proprietary software components under applicable FDA procedures;
\item[(iii)] The ability to observe Physical AI system operations in a non-clinical demonstration setting, including execution of test procedures, emergency stop activation, and fault recovery protocols; and
\item[(iv)] Digital access to hash-chained audit trails and cryptographically signed log files, with documentation of the verification procedures for confirming data integrity.
\end{description}

[52 FR 8831, Mar. 19, 1987, as amended at 67 FR 9586, Mar. 4, 2002. Physical AI adaptation added 17 March 2026.]

\subsection{\S~312.59 Disposition of Unused Supply of Investigational Drug}

The sponsor shall assure the return of all unused supplies of the investigational drug from each individual investigator whose participation in the investigation is discontinued or terminated. The sponsor may authorize alternative disposition of unused supplies of the investigational drug provided this alternative disposition does not expose humans to risks from the drug.

\subsubsection{Physical AI Adaptation of \S~312.59}

When the participation of an investigator is discontinued or terminated and the investigation involved Physical AI systems:

\begin{description}[leftmargin=0.8cm,style=sameline]
\item[(a)] The sponsor shall ensure that all Physical AI systems at the investigator's site are placed in safe standby mode and are decommissioned or retrieved in accordance with the decommissioning plan specified in the Physical AI System Description;
\item[(b)] All Physical AI system data, including telemetry archives, system logs, and configuration data, shall be securely transferred to the sponsor or securely destroyed in accordance with the data management plan;
\item[(c)] All Physical AI system access credentials, authorization tokens, and MCP server configurations at the site shall be revoked; and
\item[(d)] The investigator shall certify in writing that all Physical AI system components have been returned, decommissioned, or disposed of in accordance with the sponsor's instructions, and that no Physical AI system at the site retains access to the investigational drug administration systems.
\end{description}

[52 FR 8831, Mar. 19, 1987, as amended at 67 FR 9586, Mar. 4, 2002. Physical AI adaptation added 17 March 2026.]

\subsection{\S~312.60 General Responsibilities of Investigators}

An investigator is responsible for ensuring that an investigation is conducted according to the signed investigator statement, the investigational plan, and applicable regulations; for protecting the rights, safety, and welfare of subjects under the investigator's care; and for the control of drugs under investigation. An investigator shall, in accordance with the provisions of part 50 of this chapter, obtain the informed consent of each human subject to whom the drug is administered, except as provided in \S~50.23 or \S~50.24. Additional specific responsibilities of clinical investigators are set forth in this part and in parts 50 and 56.

\subsubsection{Physical AI Adaptation of \S~312.60}

In addition to the general responsibilities described above, when an investigation involves Physical AI systems, each investigator shall be responsible for:

\begin{description}[leftmargin=0.8cm,style=sameline]
\item[(a)] \textit{Physical AI system operation.} Ensuring that the Physical AI system is operated in accordance with the approved protocol and the Physical AI System Description, including adherence to the authorized procedure types, drug administration parameters, and safety limits;
\item[(b)] \textit{Human oversight.} Maintaining continuous human oversight of the Physical AI system during all clinical procedures involving investigational drugs, including the ability to immediately assume manual control or activate the emergency stop at any time;
\item[(c)] \textit{Pre-procedure verification.} Completing the pre-procedure safety matrix verification before each Physical AI-assisted clinical procedure, confirming system readiness, patient-specific parameter configuration, drug preparation verification, and environmental safety checks;
\item[(d)] \textit{Physical AI adverse event documentation.} Documenting all Physical AI adverse events that occur during clinical procedures under the investigator's supervision, including the event classification, system state at the time of the event, operator actions taken, and patient outcome;
\item[(e)] \textit{Physical AI system integrity.} Reporting to the sponsor any observed Physical AI system anomalies, degraded performance, or unexpected behaviors, even when such observations do not rise to the level of a reportable Physical AI adverse event;
\item[(f)] \textit{Informed consent for Physical AI.} Ensuring that the informed consent process includes a clear explanation of the Physical AI system's role in the investigation, consistent with the requirements of \S~312.23(a)(1)(iv), including a description of the Physical AI system, its function in the clinical procedure, the human oversight protocols in place, the risks associated with the Physical AI system, and the subject's right to request that the procedure be performed without Physical AI assistance if clinically feasible; and
\item[(g)] \textit{Physical AI training currency.} Maintaining current training and competency on the Physical AI system, including completing any refresher training or requalification assessments required by the sponsor following Physical AI system updates.
\end{description}

[52 FR 8831, Mar. 19, 1987. Physical AI adaptation added 17 March 2026.]

\subsection{\S~312.61 Control of the Investigational Drug}

An investigator shall administer the drug only to subjects under the investigator's personal supervision or under the supervision of a subinvestigator responsible to the investigator. The investigator shall not supply the drug to any person not authorized under this part to receive it.

\subsubsection{Physical AI Adaptation of \S~312.61}

When the investigational drug is administered through or with the assistance of a Physical AI system:

\begin{description}[leftmargin=0.8cm,style=sameline]
\item[(a)] The Physical AI system shall be operated only by the investigator or a qualified subinvestigator who has completed the required Physical AI system training and has been authorized by the sponsor;
\item[(b)] The Physical AI system shall not be configured to administer the investigational drug autonomously without the continuous presence and supervision of a qualified investigator or subinvestigator;
\item[(c)] Access controls on the Physical AI system shall be configured to prevent unauthorized persons from initiating drug administration procedures, modifying drug administration parameters, or overriding safety limits; and
\item[(d)] The investigator shall maintain accountability records for the investigational drug that include documentation of any Physical AI-assisted administration, including the date, time, dose, route, Physical AI system identifier, and the identity of the supervising investigator or subinvestigator.
\end{description}

[52 FR 8831, Mar. 19, 1987. Physical AI adaptation added 17 March 2026.]

\subsection{\S~312.62 Investigator Recordkeeping and Record Retention}

(a) \textit{Disposition of drug.} An investigator is required to maintain adequate records of the disposition of the drug, including dates, quantity, and use by subjects. If the investigation is terminated, suspended, discontinued, or completed, the investigator shall return the unused supplies of the drug to the sponsor, or otherwise provide for disposition of the unused supplies of the drug under \S~312.59.

(b) \textit{Case histories.} An investigator is required to prepare and maintain adequate and accurate case histories that record all observations and other data pertinent to the investigation on each individual administered the investigational drug or employed as a control in the investigation. Case histories include the case report forms and supporting data including, for example, signed and dated consent forms and medical records including, for example, progress notes of the physician, the individual's hospital chart(s), and the nurses' notes.

\subsubsection{Physical AI Adaptation of \S~312.62}

(c) \textit{Physical AI procedure records.} In addition to the records required under paragraphs (a) and (b), the investigator shall maintain the following records for each Physical AI-assisted clinical procedure:

\begin{description}[leftmargin=0.8cm,style=sameline]
\item[(i)] The pre-procedure safety matrix verification record, documenting the completion and results of all pre-procedure checks;
\item[(ii)] The Physical AI system identifier (model, serial number, software version, AI/ML model version) used in the procedure;
\item[(iii)] A procedure event log documenting the Physical AI system's actions during the procedure, including drug administration events, operator interventions, safety limit activations, and any system anomalies;
\item[(iv)] Any Physical AI adverse events that occurred during the procedure, documented in accordance with the Physical AI adverse event classification system specified in \S~312.32(g); and
\item[(v)] Post-procedure system status verification, confirming that the Physical AI system is in a safe state and that all procedure data have been captured and stored.
\end{description}

(d) \textit{Physical AI case history integration.} The case histories required under paragraph (b) shall integrate Physical AI system data with clinical observations, including:

\begin{description}[leftmargin=0.8cm,style=sameline]
\item[(i)] Notation that the procedure was Physical AI-assisted and the specific Physical AI system used;
\item[(ii)] Any deviations from the planned Physical AI procedure, including instances where the investigator assumed manual control, activated the emergency stop, or modified Physical AI system parameters during the procedure; and
\item[(iii)] The investigator's clinical assessment of the Physical AI system's performance during the procedure, including any observations relevant to patient safety.
\end{description}

[52 FR 8831, Mar. 19, 1987, as amended at 52 FR 23031, June 17, 1987; 61 FR 57280, Nov. 5, 1996; 67 FR 9586, Mar. 4, 2002. Physical AI adaptation added 17 March 2026.]

\subsection{\S~312.64 Investigator Reports}

(a) \textit{Progress reports.} The investigator shall furnish all reports to the sponsor of the drug who is responsible for collecting and evaluating the results obtained.

(b) \textit{Safety reports.} An investigator shall promptly report to the sponsor any adverse effect that may reasonably be regarded as caused by, or probably caused by, the drug. If the adverse effect is alarming, the investigator shall report the adverse effect immediately.

(c) \textit{Final report.} An investigator shall provide the sponsor with an adequate report shortly after completion of the investigator's participation in the investigation.

\subsubsection{Physical AI Adaptation of \S~312.64}

(d) \textit{Physical AI progress reports.} In addition to the progress reports required under paragraph (a), the investigator shall furnish to the sponsor periodic reports on Physical AI system performance at the investigator's site, including aggregate procedure counts, Physical AI adverse event summaries, system uptime and reliability metrics, and any training or requalification activities completed during the reporting period.

(e) \textit{Physical AI safety reports.} In addition to the safety reports required under paragraph (b):

\begin{description}[leftmargin=0.8cm,style=sameline]
\item[(i)] The investigator shall promptly report to the sponsor any Physical AI adverse event, as defined in \S~312.32(g), regardless of whether the event resulted in patient harm;
\item[(ii)] Physical AI adverse events that are alarming, including uncontrolled robotic movement, drug administration at incorrect doses by the Physical AI system, or loss of emergency stop functionality, shall be reported immediately to the sponsor; and
\item[(iii)] The investigator shall report to the sponsor any observed Physical AI system behavior that, in the investigator's clinical judgment, could compromise patient safety if the behavior were to recur or escalate, even if the observation does not meet the formal definition of a Physical AI adverse event.
\end{description}

(f) \textit{Physical AI final report.} The final report required under paragraph (c) shall include a summary of the investigator's experience with the Physical AI system, including an assessment of the system's reliability, usability, and clinical utility, and any recommendations for Physical AI system improvements based on the clinical investigation experience.

[52 FR 8831, Mar. 19, 1987, as amended at 52 FR 23031, June 17, 1987; 67 FR 9586, Mar. 4, 2002. Physical AI adaptation added 17 March 2026.]

\subsection{\S~312.66 Assurance of IRB Review}

An investigator shall assure that an IRB that complies with the requirements set forth in part 56 will be responsible for the initial and continuing review and approval of the proposed clinical study. The investigator shall also assure that he or she will promptly report to the IRB all changes in the research activity and all unanticipated problems involving risk to human subjects or others, and that he or she will not make any changes in the research without IRB approval, except where necessary to eliminate apparent immediate hazards to human subjects.

\subsubsection{Physical AI Adaptation of \S~312.66}

When a clinical investigation involves Physical AI systems, the investigator shall additionally assure that:

\begin{description}[leftmargin=0.8cm,style=sameline]
\item[(a)] The IRB has been provided with sufficient information about the Physical AI system to conduct a meaningful review of the risks and benefits, including the Physical AI System Description or a summary thereof, the informed consent language regarding the Physical AI system, and a description of the human oversight protocols;
\item[(b)] All changes to the Physical AI system that may affect the risk profile of the investigation, including software updates, AI/ML model changes, and hardware modifications, are reported to the IRB in accordance with its continuing review requirements;
\item[(c)] All Physical AI adverse events that constitute unanticipated problems involving risk to human subjects are promptly reported to the IRB; and
\item[(d)] The IRB is notified of any Physical AI system-related clinical holds, FDA communications regarding Physical AI system deficiencies, or Physical AI-specific amendments to the protocol.
\end{description}

[52 FR 8831, Mar. 19, 1987, as amended at 67 FR 9586, Mar. 4, 2002. Physical AI adaptation added 17 March 2026.]

\subsection{\S~312.68 Inspection of Investigator's Records and Reports}

An investigator shall upon request from any properly authorized officer or employee of FDA, at reasonable times, permit such officer or employee to have access to, and copy and verify any records or reports made by the investigator pursuant to \S~312.62. The investigator is not required to divulge subject names unless the records of particular individuals require a more detailed study of the cases, or unless there is reason to believe that the records do not represent actual cases studied or results obtained.

\subsubsection{Physical AI Adaptation of \S~312.68}

In addition to the records and reports described above, the investigator shall make available to FDA for inspection all Physical AI-specific records required under \S~312.62(c) and (d), including Physical AI procedure event logs, pre-procedure safety matrix records, Physical AI adverse event documentation, and case histories integrating Physical AI system data. Upon FDA request, the investigator shall also facilitate access to the Physical AI system for the purpose of a non-clinical demonstration of system capabilities, emergency stop functionality, and safety protocols.

[52 FR 8831, Mar. 19, 1987, as amended at 67 FR 9586, Mar. 4, 2002. Physical AI adaptation added 17 March 2026.]

\subsection{\S~312.69 Handling of Controlled Substances}

If the investigational drug is subject to the Controlled Substances Act, the investigator shall take adequate precautions, including storage of the investigational drug in a securely locked, substantially constructed cabinet, or other securely locked, substantially constructed enclosure, access to which is limited, to prevent theft or diversion of the substance into illegal channels of distribution.

\subsubsection{Physical AI Adaptation of \S~312.69}

When a Physical AI system is used to administer an investigational drug that is a controlled substance, the following additional precautions apply:

\begin{description}[leftmargin=0.8cm,style=sameline]
\item[(a)] The Physical AI system's drug storage and handling compartments shall be secured with physical locks and electronic access controls that limit access to authorized investigators and subinvestigators only;
\item[(b)] The Physical AI system shall maintain a tamper-evident, hash-chained log of all controlled substance loading, administration, and disposal events, including the identity of the authorizing investigator, the quantity loaded and administered, and any discrepancies; and
\item[(c)] The Physical AI system's drug administration pathway shall be designed to prevent unauthorized diversion, including physical safeguards against unauthorized removal of the drug from the system and software controls that reconcile drug quantities loaded against quantities administered and returned.
\end{description}

[52 FR 8831, Mar. 19, 1987, as amended at 67 FR 9586, Mar. 4, 2002. Physical AI adaptation added 17 March 2026.]

\subsection{\S~312.70 Disqualification of a Clinical Investigator}

(a) If FDA has information indicating that an investigator (including a sponsor-investigator) has repeatedly or deliberately failed to comply with the requirements of this part, part 50, or part 56, or has submitted to the sponsor or FDA misleading reports, FDA will furnish the investigator written notice of the matter, offer an opportunity for explanation, and offer an opportunity for a regulatory hearing under part 16 on the question of whether the investigator is entitled to receive investigational drugs.

(b) On the basis of the administrative record or a regulatory hearing, if the Commissioner determines that the investigator has repeatedly or deliberately violated the requirements, the Commissioner may prohibit the investigator from receiving investigational drugs. Such notice will include a statement of the basis for such determination and may be revoked if adequate assurances of future compliance are provided.

(c) Each IND and each approved application shall include the name of each investigator who has been determined to be ineligible, together with a brief statement of the basis for such determination. This list shall be updated within 15 working days of such determination.

\subsubsection{Physical AI Adaptation of \S~312.70}

(d) \textit{Physical AI grounds for disqualification.} In addition to the grounds specified in paragraph (a), FDA may furnish written notice to an investigator who has:

\begin{description}[leftmargin=0.8cm,style=sameline]
\item[(i)] Repeatedly or deliberately operated a Physical AI system outside the parameters authorized by the approved protocol and Physical AI System Description;
\item[(ii)] Failed to maintain required human oversight of Physical AI system operations during clinical procedures, including operating the Physical AI system without the required qualified supervision;
\item[(iii)] Failed to report Physical AI adverse events to the sponsor as required by this part;
\item[(iv)] Submitted misleading Physical AI system performance data, training records, or Physical AI adverse event reports to the sponsor or to FDA; or
\item[(v)] Allowed unauthorized personnel to operate or access the Physical AI system in connection with the clinical investigation.
\end{description}

(e) \textit{Physical AI competency suspension.} When an investigator is determined to be ineligible under this section based in whole or in part on Physical AI-related grounds, the determination shall specify whether the investigator is prohibited from:

\begin{description}[leftmargin=0.8cm,style=sameline]
\item[(i)] Receiving investigational drugs generally;
\item[(ii)] Receiving investigational drugs only in connection with Physical AI-assisted investigations; or
\item[(iii)] Operating or supervising any Physical AI system in a clinical investigation, while remaining eligible to conduct non-Physical AI investigations.
\end{description}

[52 FR 8831, Mar. 19, 1987, as amended at 52 FR 23031, June 17, 1987; 67 FR 9586, Mar. 4, 2002. Physical AI adaptation added 17 March 2026.]
